\documentclass[11pt,a4paper]{ctexart}
\usepackage[margin=2.1cm]{geometry}
\usepackage{amsmath,amssymb,booktabs,float,graphicx,xcolor,adjustbox,longtable,hyperref,enumitem,array}
\usepackage{placeins}
\hypersetup{colorlinks=true,linkcolor=blue!50!black,urlcolor=blue!50!black}
\setlist{itemsep=2pt,topsep=3pt}
\title{高对比度问题的求解器：\\目前的方法、适用场景和实验数据}
\author{温州肯恩大学 \quad Nonlinear Poisson 项目}
\date{2026 年 11 月 11 日}

\begin{document}
\maketitle
\sloppy

\begin{abstract}
本文汇总 11 月 3 日到 11 月 9 日为高对比度问题做的四个求解器：FD-LGF、bt-fd、FluxMG 和 SSR，写清每个方法的原理、适合的场景和已有的实验数据，最后给出我目前的判断和下一步。附录 A--E 用表格列出全部实验数据。本文没有新实验，数据全部取自已有的结果文件（见第~\ref{sec:files} 节）。

\textbf{主要判断。}(1) 标量系数、通量差分的高对比度问题（M 矩阵）上，BoomerAMG 的迭代数很稳定，$\chi$ 到 $10^8$ 都只要 7--13 次。要胜过它，只能靠利用问题结构降低每步和构建的成本，不能指望更好的鲁棒性。
(2) 没有一个方法在所有场景都最好，所以打算做成按矩阵结构选择的组合：背景可分离用 FD-LGF（三维 $\chi^x$ 上约为 BoomerAMG 的 0.04 倍）；三温问题且特征区小用 bt-fd（SC26\_01 上约 0.5 倍）；张量网格上的多场和极端对比用 FluxMG（留出集对 BoomerAMG 最优设置 9/10 获胜，真实矩阵 0.80 倍）；张量网格单场用 hypre PFMG（比 FluxMG 快 1.8--3.4 倍）；其余情况仍用 BoomerAMG。
(3) 这套选择规则只有部分阈值有数据支持，还没有作为整体验证。所有结果都是单线程，除 SC26\_01 外都是我们自己生成的矩阵。
\end{abstract}

\tableofcontents

\section{“高对比度”在本项目中的三种含义}
\begin{enumerate}
\item \textbf{系数跨度大。}$-\nabla\cdot(a\nabla u)+qu=f$，教授指定的主算例是 $a=a_0\chi^x$，对比度恰好是 $\chi$；补充算例有夹杂、通道、绝缘墙、对数正态场、渗流两相等，统一写成 $a=\chi^{g}$，$g\in[0,1]$，$\chi\le10^8$。
\item \textbf{极端量级加多场耦合。}三温辐射扩散（3T），如 SC26\_01：离子导热跨过冷密壳层下降约 $10^{19.6}$ 倍，壳层内三个温度之间的交换系数比别处大 $10^3$--$10^{19}$ 倍，矩阵中有 $10^{17}$ 量级的元素。普通的“对角元减非对角元之和”会严重相消，热核形成近乎奇异的孤岛。
\item \textbf{多群耦合。}多群扩散的群间耦合强弱差别很大，选错哪些场一起解，迭代数可以从 23 变成 121。我们的多群算例只是类比 Yue 等 2026 的难易程度，不是他们的矩阵。
\end{enumerate}
三类问题上 BoomerAMG 的弱点不同：第 1 类上迭代数稳定，弱点是三维构建成本（$n=127$ 时 setup 5--8\,s）；第 2、3 类上，它把多场当成一个标量矩阵处理，迭代数和成本都更高，难的多群算例要 100--300 次。

\section{一页结论：场景与推荐方法}\label{sec:choice}
\begin{table}[H]
\centering\small
\renewcommand{\arraystretch}{1.25}
\begin{tabular}{>{\raggedright}p{4.3cm}>{\raggedright}p{2.2cm}>{\raggedright}p{4.4cm}>{\raggedright\arraybackslash}p{4.3cm}}
\toprule
场景 & 推荐 & 实测（对 BoomerAMG） & 不适用 / 尚缺证据 \\
\midrule
笛卡尔网格，背景可分离或分层（$\chi^x$、阶跃层、薄层），特征体积小 & FD-LGF & 三维 $\chi^x$，$n=127$：0.04 倍，迭代 2--3 次；三维薄层 25--40 倍更快 & 不可分离系数退化为直接法，慢 1.8--14 倍；未与 PFMG 比较 \\
三温，单一背景占主导，特征区 $\lesssim10\%$ & bt-fd（Cholesky） & SC26\_01：0.51 倍；特征区 4.7--9.1\% 的序列 0.53--0.73 倍（对调参后的 BoomerAMG） & 特征区 16--20\% 时 1.03--1.35 倍，输；并行未测 \\
张量网格 7 点，多场（3T 多夹杂、多群），极端对比 & FluxMG & 留出集 9/10 获胜，SCORE $+0.65$；SC26\_01 0.80 倍（对 BoomerAMG 最优设置） & 多群易例 1.59 倍；对齐各向异性 1.15--1.22 倍；非张量网格不支持 \\
张量网格，单场，一般高对比度（随机场、管道、绝缘板、渗流） & hypre PFMG（红黑） & 比 FluxMG 快 1.9--3.0 倍，约为 BoomerAMG 的 0.3--0.4 倍 & 3T 胶囊类和 SC26\_01 上失败 \\
非结构网格、一般稀疏矩阵 & BoomerAMG & -- & SSR 在高对比度标量上 0.95--1.18 倍，不推荐 \\
\bottomrule
\end{tabular}
\caption{各场景的推荐方法。比值小于 1 表示比 BoomerAMG 快。不同行的数据来自不同实验，判据不同，见各节。}
\end{table}

\section{FD-LGF：可分离背景的精确逆加局部精确解}
\subsection{方法}
FD-LGF 是 Fictitious-Domain Lattice Green's Function 预条件。把系数分解为
\[
  a(x,y)=b(x,y)\,\phi(x,y),\qquad b=a_1(x)\,a_2(y),\qquad \phi\ne1\ \text{只在特征（夹杂、通道、墙）内}.
\]
\begin{itemize}
\item \textbf{背景逆 $G=B^{-1}$。}$b$ 可分离时，通量算子可写成 $BU=D_2UT_x+T_yUD_1$，用快速对角化（Lynch--Rice--Thomas 1964）精确求逆。$a_2$ 为常数时 $y$ 方向用 DST-I，复杂度 $O(n^2\log n)$；两个方向都变化时用稠密形式，二维每次 $O(n^3)$。
\item \textbf{局部精确解 $Q$。}特征集 $\{|\log a-\log b|>\log2\}$ 外扩宽度 $w$ 的环，按连通分量分块做稀疏 LU，$Q=\sum_kR_k^{\mathsf T}A_k^{-1}R_k$。必须精确：换成 V-cycle 或 Chebyshev 后迭代升到 25--540 次。
\item \textbf{平衡组合。}$P=Q+(I-QA)\,G\,(I-AQ)$，对称正定，直接用于 PCG。
\end{itemize}
$G$ 对背景精确、$A_k^{-1}$ 对局部精确，两者都不随 $\chi$ 变差，所以迭代数与 $\chi$ 无关；环宽按物理长度固定时也与 $n$ 无关。

\subsection{适用场景}
笛卡尔网格，背景（近似）可分离或分层，特征体积小。典型例子是教授的 $\chi^x$、分层储层加少量夹杂或断层、三维层状复合材料。特征体积占大部分时，$Q$ 就是一次整体 LU，退化为稀疏直接法，不应使用。

\subsection{实验数据}
判据：PCG，$\|r\|/\|f\|\le10^{-10}$，单线程，3 次取中位数；BoomerAMG 为 PETSc 默认参数（三维另测 $\theta=0.25,0.5$）。离散 MMS，所有收敛方法得到同一个离散解。

\begin{table}[H]
\centering\small
\begin{tabular}{llrrrr}
\toprule
算例（$\chi=10^8$） & $n$ & 传统 LGF & GMG & BoomerAMG & FD-LGF \\
\midrule
二维 A：$\chi^x$ & 1023 & 46.30 & 0.90 & 1.14 & \textbf{0.25} \\
二维 B：阶跃层 & 1023 & 3.04 & 0.94 & 1.35 & \textbf{0.20} \\
二维 C：$\chi^x$+规则夹杂 & 1023 & 38.69 & 8.20 & \textbf{1.53} & 2.26 \\
二维 D：$\chi^x$+随机夹杂 & 1023 & 37.92 & 9.33 & \textbf{1.62} & 2.72 \\
\midrule
三维 A：$\chi^x$ & 127 & 47.96 & 2.62 & 8.86 & \textbf{0.34} \\
三维 B：阶跃层 & 127 & 5.14 & 2.28 & 12.80 & \textbf{0.23} \\
三维 C：$\chi^x$+规则夹杂 & 127 & 44.08 & 13.50 & 10.09 & \textbf{9.25} \\
三维 D：$\chi^x$+随机夹杂 & 127 & 43.00 & 19.73 & 11.66 & \textbf{6.01} \\
\bottomrule
\end{tabular}
\caption{教授的 $\chi^x$ 系列，总时间（s，setup+solve）。夹杂算例的 FD-LGF 取环宽 $w=1/64$ 的结果。迭代数：A、B 上 FD-LGF 1--3 次，与 $\chi$、$n$ 无关；传统 LGF 在 $\chi\ge10^4$ 时不收敛；三维 BoomerAMG 在夹杂算例上随 $n$ 从 11 增到 17 次，FD-LGF 为 10--16 次。数据：\texttt{FD\_LGF/report/tables\_2d.tex, tables\_3d.tex}。}
\end{table}

\begin{figure}[H]\centering
\includegraphics[width=0.62\textwidth]{fig_fdlgf.pdf}
\caption{三维 $\chi^x$（$\chi=10^8$）的总时间随网格的变化。FD-LGF 在 $n=127$（约 $2\times10^6$ 未知量）时为 BoomerAMG 的 0.04 倍。}
\end{figure}

\begin{table}[H]
\centering\small
\footnotesize\setlength{\tabcolsep}{4pt}
\begin{tabular}{lrrrl}
\toprule
系数（二维 $n=1023$，4 个 $\chi$） & AMG 迭代 & FD-LGF 迭代 & 时间比 & 局部区域比例 \\
\midrule
sep2 两方向可分离 & 8--9 & 3--4 & 0.9--1.1 & 0 \\
walls 绝缘墙 & 10--11 & 9--14 & 1.1--1.2 & 0.24--0.64 \\
chan 通道 & 12--13 & 9--25 & 1.4--2.1 & 0.26--0.65 \\
\begin{tabular}[t]{@{}l@{}}其余 7 类（checker、diag、layered、\\ lognorm、perc、diagfine、rough）\end{tabular} & 8--13 & 1--28 & 1.8--14 & 0.20--1.00 \\
\midrule
三维 sheets 薄层（$n=127$） & 9--11 & 2--3 & 0.02--0.04 & -- \\
三维 tubes 管道（$n=127$，$w=1/64$） & 10--12 & 7--17 & 约 0.5 & -- \\
三维 insul 绝缘板（$n=127$） & 8--12 & 8--9 & 9--100 & 贯穿全域 \\
\midrule
\multicolumn{5}{l}{时间比 = FD-LGF / BoomerAMG 总时间}\\
\bottomrule
\end{tabular}
\caption{更一般的 10 类二维系数和 5 类三维系数（报告 01 的确认实验，方案锁定 \texttt{c2ff7672}）。BoomerAMG 在所有情形下都没有迭代失败。数据：\texttt{FD\_LGF/high\_contrast/report\_hc/}。}
\end{table}

\paragraph{精度提醒。}按残差停止时，FD-LGF 在低系数区的代数误差比 BoomerAMG 大 1--4 个数量级（棋盘上可到 $10^{-4}$），残差判据看不到低系数区的误差。比较时要么收紧容差，要么按误差计数。

\section{bt-fd：三温问题的块三角加 FD-LGF}
\subsection{方法}
针对 SC26\_01 这类三温矩阵：在真实解上，$T_1$ 对 $T_2$ 方程的反馈只占 $10^{-12}$，电子块的非对角部分只占 $2\times10^{-8}$，所以 $3\times3$ 系统可以化为“两个标量扩散方程加一次对角消元”。
\begin{itemize}
\item \textbf{bt}：消去电子方程，按块三角顺序解剩下的两个标量块；
\item \textbf{fd}：每个标量块用 FD-LGF 的思路，常系数背景用快速变换，冷密壳层等特征区用直接解（LU 或 Cholesky）；
\item 外层用双双精度残差做混合精度迭代修正，2--3 步达到 $\|b-Ax\|/\|b\|<10^{-8}$。
\end{itemize}

\subsection{适用场景}
三温或类似的少数场耦合问题，单一常系数背景占主导，特征区体积小。文献调研（报告 25）的结论是“bt”部分是三温问题的主流结构，应保留；“fd”部分只适合这类矩阵，不作为通用方法。

\subsection{实验数据}
判据：双双精度原残差 $<10^{-8}$，时间含构建，单线程，独立进程。方案锁定（H57，sha256 \texttt{82dbdb84}），76 次运行全部达标。
\begin{table}[H]
\centering\small
\footnotesize\setlength{\tabcolsep}{4pt}
\begin{tabular}{lrrrrr}
\toprule
& 特征区 & BoomerAMG 调参 & bt-v（内含 AMG） & bt-fd Cholesky & Cholesky/调参 \\
\midrule
真实 SC26\_01（5 轮中位数） & -- & 31.86 & 29.80 & \textbf{16.78} & 0.53 \\
G1 球（10 步） & 4.7\% & 220.2 & 208.7 & \textbf{117.2} & 0.53 \\
G4 长椭球 & 9.1\% & 219.0 & 199.0 & \textbf{141.8} & 0.65 \\
G7 3 胶囊 & 6.6\% & 215.0 & 225.7 & \textbf{156.8} & 0.73 \\
G9 10 胶囊 & 16.3\% & 204.7 & \textbf{189.6} & 210.2 & 1.03 \\
G8 10 胶囊 & 18.4\% & 209.7 & \textbf{194.0} & 283.5 & 1.35 \\
G10 10 扁方块 & 20.3\% & 207.8 & \textbf{188.4} & 257.7 & 1.24 \\
\bottomrule
\end{tabular}
\caption{时间（s）。G 序列为 $m=128$、每个 10 个时间步的自建三温序列（真实矩阵只有一个时间步）。对未调参的基线 BoomerAMG，SC26\_01 上 bt-fd Cholesky 的配对比值为 0.51。数据：\texttt{FD\_LGF/regions/report/report\_h57.tex}。}
\end{table}
特征区在 4.7--9.1\% 时 bt-fd 胜，16.3--20.3\% 时输，中间没有数据，所以“$\lesssim10\%$”这个阈值只是暂定。

\section{FluxMG：通量形式的结构多重网格}
\subsection{方法}
\begin{itemize}
\item \textbf{算子只存通量。}每个场每个单元存面传导系数 $k$ 和行余量 $C$（双双精度累加得到），不存对角元：
$(Ax)_{f,i}=\sum_g C_{fg,i}\,x_{g,i}+\sum_{\text{面}}k_f(x_{f,i}-x_{f,\mathrm{nb}})$。先求差再乘，$10^{17}$ 的元素不会相消。
\item \textbf{层次。}$2\times2\times2$ 分片常数聚合，粗层面系数是跨盒子边界的细面系数之和，每层仍是 7 点格式；某方向明显偏弱时不粗化该方向（半粗化）。构建基本只有加法。
\item \textbf{循环。}K-cycle（Notay，两步 GCR），块红黑光滑，只在连续 4 个以上各向异性单元的线段上做线消元（GTH 形式，只用加法）。
\item \textbf{区域向量。}每层把扩散占优的相连单元分成区域，每个区域一个常数向量，经静态消元延伸到相邻单元，小矩阵 $E=W^{\mathsf T}AW$ 按差分精确计算后用 GTH 消元。这一步专门处理孤岛和近乎奇异的高导区域。
\item \textbf{多场缩减。}可由对角解的场先消去，其余按常数模态的回路增益分组，块三角顺序求解，每组一个 Core 循环；再用一个探针：在聚合到 $8^3$ 盒子的系统上测块三角求解的收缩因子，超过 0.7 就把所有耦合场并成一组。
\end{itemize}
全程不用 AMG，只需要张量网格上的 7 点 M 矩阵系统。

\subsection{适用场景}
张量网格上的多场问题（三温多夹杂、多群），以及系数量级极端、普通存储会相消的问题。单场问题上它能胜过 BoomerAMG，但这个优势来自结构网格本身，hypre 的结构多重网格 PFMG 更快（见下）。

\subsection{实验数据}
判据：双双精度原残差 $<10^{-8}$，时间含构建，单线程。对照组有两个：BoomerAMG 在 8 组标量设置加 sys.u/sys.n（多场的 systems 设置）中每个算例取最快的一个；hypre PFMG、PFMG 红黑、SMG 中取最快的一个（预注册比较，2 次打乱顺序的重复）。
\begin{figure}[H]\centering
\includegraphics[width=0.86\textwidth]{fig_fluxmg.pdf}
\caption{测试集 19 个算例加真实矩阵，时间对 BoomerAMG 最优设置的比值。上面 7 个是单场（Poisson、各向异性、高对比度），中间 9 个是三温，再下面 3 个是多群，最下面是真实矩阵。数据：\texttt{fluxmg/structured/results\_20261109.jsonl}、\texttt{fluxmg/amg\_configs.jsonl}。}
\end{figure}

\begin{table}[H]
\centering\small
\begin{tabular}{lrrl}
\toprule
算例集 & 对 BoomerAMG 最优设置 & 对 PFMG/SMG 最优 & 说明 \\
\midrule
测试集（19） & $+0.43$，14/19 胜 & $+0.22$，8/19 胜 & 循环中使用过 \\
留出集（10） & $+0.65$，9/10 胜 & $+1.72$，9/10 胜 & 循环中从未运行 \\
多群分组留出（14） & $+0.57$，11/14 胜 & $+0.60$，9/14 胜 & 分组规则实现前固定 \\
真实 SC26\_01 & 0.80 倍 & PFMG/SMG 全部失败 & 59 次 vs 34 次 \\
\bottomrule
\end{tabular}
\caption{SCORE 为 $-\mathrm{mean}\log_2(t_\text{FluxMG}/t_\text{对照})$，越大越好，失败按 8 倍计；$+0.43$ 约相当于平均快 1.35 倍。}
\end{table}

\paragraph{各类问题上的情况。}
\begin{itemize}
\item \textbf{单场高对比度}（lognorm3、tubes、insul、perc3，$64^3$）：FluxMG 0.39--0.72\,s，BoomerAMG 最优设置 0.60--0.71\,s，PFMG 红黑 0.20--0.24\,s。FluxMG 对 BoomerAMG 为 0.65--1.02 倍，但 PFMG 又比 FluxMG 快 1.9--3.0 倍。
\item \textbf{三温}：FluxMG 对 BoomerAMG 为 0.47--0.86 倍（10 渐变扁方块 1.07 倍是唯一输的）；PFMG/SMG 在胶囊类、$m=128$ 和真实矩阵上全部失败，这是 FluxMG 区域向量起作用的地方。
\item \textbf{多群}：难例 0.29--0.53 倍；易例 1.59 倍（对 sys.n），中例 1.04 倍。
\item \textbf{对齐各向异性}（强方向 100、1000）：1.15--1.22 倍。
\item 真实矩阵上 FluxMG 约 23\,s，bt-fd Cholesky 约 16.8\,s（另一批运行，判据相同），所以在小特征区的三温问题上 bt-fd 仍然更好。
\end{itemize}

\paragraph{成本。}构建是主要开销，其中 \texttt{build\_regions}/\texttt{assemble\_regions} 约占 60\%；但区域向量只对三温有用，而且三温在最细层和粗层都需要。试过用粗层反应项的非 Galerkin 占比、单区域判据、探针来决定是否建区域向量，都分不开（第 31 轮）。预先算好 $H=AW$、$G=A^{\mathsf T}W$ 让 Poisson 的预条件从 8.6 降到 6.2\,ms，测试集上没有收益（第 33 轮）。

\section{SSR：纯代数的采样 Schur 约化（作为对照）}
SSR 只读 CSR 矩阵：每层对独立集精确消元，用随机采样（对称时链采样，非对称时循环运输方案）代替 Schur 补的填充，再做 GTH 余量更新和同号合并。它是为“比 BoomerAMG 更通用”设计的，不是专门的高对比度方法，放在这里是为了说明纯代数路线在高对比度上的表现。
\begin{table}[H]
\centering\small
\begin{tabular}{lrrr}
\toprule
算例 & SSR 迭代/时间 & BoomerAMG \texttt{fast.25} & SSR/AMG \\
\midrule
lognorm3 & 25 / 0.66\,s & 19 / 0.56\,s & 1.18 \\
tubes & 35 / 0.78\,s & 20 / 0.66\,s & 1.18 \\
perc3 & 27 / 0.65\,s & 19 / 0.68\,s & 0.95 \\
insul（留出） & 26 / 0.65\,s & 19 / 0.60\,s & 1.09 \\
3T（3 例） & 25--46 & 30 & 1.12--1.64 \\
多群难例 & 22 / 0.47\,s & 109 / 2.36\,s & 0.20 \\
\bottomrule
\end{tabular}
\caption{锁定版本 G（\texttt{gensolver/round3/final\_dev\_G.log}、\texttt{final\_holdout\_G.log}）。}
\end{table}
高对比度标量和三温上 SSR 略慢于 BoomerAMG，只在多群、图和对流类问题上明显更快。另外一个教训：按 $\sqrt{d_id_j}$ 判据丢弃小元素会破坏高对比度界面，这条路已放弃。

\section{我目前的判断和下一步}\label{sec:next}
\subsection{判断}
\begin{enumerate}
\item \textbf{要赢靠结构，不靠鲁棒性。}在 M 矩阵的高对比度问题上，BoomerAMG 迭代数从来不是问题（报告 01 没有找到它迭代失败的例子）。我们几个方法的收益都来自利用结构：FD-LGF 利用可分离背景，bt-fd 利用三温的弱反馈，FluxMG 和 PFMG 利用张量网格。
\item \textbf{应当做组合，而且选择规则必须只看矩阵。}可以从矩阵算出的判别量有：场的个数和耦合强度（FluxMG 的分组探针已经实现）、是否张量网格 7 点格式、背景的可分离程度（对 $\log a$ 做可分离拟合的残差）、特征区体积分数。暂定的规则就是第~\ref{sec:choice} 节的表。
\item \textbf{FluxMG 的价值在多场和极端对比，不在单场。}11 月 9 日的预注册比较里预测 P5（“FluxMG 在 70\% 以上的测试和留出算例上快于最好的结构多重网格”）不成立（59\%）。单场问题应直接交给 PFMG。
\item \textbf{高对比度要看真实误差。}残差判据会掩盖低系数区的误差（FD-LGF 上差 1--4 个数量级），三温上的自报残差也不等于真实精度。以后的比较统一用双双精度原残差，必要时再加误差判据。
\end{enumerate}

\subsection{正在做和准备做的事（均未验证）}
\begin{enumerate}
\item \textbf{FluxMG 的算子相关插值。}工作目录里的 \texttt{flux\_mg.py} 已加入几种选项：沿粗化方向的算子相关线性插值（串联电阻形式）、转置限制、粗层面系数按粗化比缩放或按串联方式计算（Kwak；Khalil--Wesseling；Alcouffe 等 1981）。这些还没有跑评测，\texttt{best/} 没有改动。目标是渐变过渡和 10 渐变扁方块（67 次）这类慢收敛算例。
\item \textbf{降低区域向量的构建成本}，或者找到一个只看矩阵就能决定是否需要区域向量的判据。
\item \textbf{弄清单场上 FluxMG 为什么比 PFMG 慢 2--3 倍。}猜测是 K-cycle 的 GCR、区域向量和仍在 Python 中的构建带来的开销，需要逐层计时。
\item \textbf{补齐对照}：FD-LGF 的 $\chi^x$ 和 10 类系数上没有和 PFMG 比过；bt-fd 和 FluxMG 之间的特征区阈值（10\%--16\% 之间）没有数据。
\item \textbf{并行测试。}全部结果都是单线程。红黑光滑和结构化存储应当利于并行，但这只是预期；“是否真正胜过 BoomerAMG”最终取决于这一步。
\item \textbf{更大网格}（$m=256$）的时间和内存，以及时间序列复用（真实矩阵只有一个时间步，只能用生成器）。
\end{enumerate}

\section{实验条件和注意事项}
\begin{itemize}
\item 全部单线程（\texttt{OMP\_NUM\_THREADS=1}），在开着其他程序的桌面机上运行；同一代码的 FluxMG SCORE 在 0.43--0.56 之间波动，所以只比较同批或配对的结果。
\item 除 SC26\_01 外，所有矩阵都由我们自己的生成器产生；多群算例是 Yue 等 2026 的类比，不是他们的矩阵；留出集和测试集来自同样的生成器，只是换了种子和参数。
\item 判据不统一：FD-LGF 系列（11 月 3 日）用 PCG 和 $\|r\|/\|f\|\le10^{-10}$；bt-fd 和 FluxMG（11 月 6--9 日）用双双精度原残差 $10^{-8}$。不同节之间的比值不能直接相比。
\item BoomerAMG 对照：FD-LGF 系列为默认参数（三维另测 $\theta$）；FluxMG 为每个算例 8--10 组设置中的最优；bt-fd 为基线和 12 组中筛出的调参版本。
\end{itemize}

\section{数据来源}\label{sec:files}
\begin{itemize}
\item FD-LGF：\texttt{FD\_LGF/report/}（$\chi^x$ 系列），\texttt{FD\_LGF/high\_contrast/report\_hc/}（10 类二维和 5 类三维系数）；报告 01。
\item bt-fd：\texttt{FD\_LGF/regions/report/report\_h57.tex}、\texttt{report\_audit\_fix.tex}、\texttt{report\_bt\_alternatives.tex}；报告 24--26。
\item FluxMG：\texttt{FD\_LGF/fluxmg/}（\texttt{final\_suite.txt}、\texttt{final\_holdout.txt}、\texttt{autoresearch.jsonl}、\texttt{amg\_configs.jsonl}、\texttt{structured/results\_20261109.jsonl}）；报告 27。
\item SSR：\texttt{FD\_LGF/gensolver/round3/}（锁定版本 G）；报告 28。
\item 图：\texttt{FD\_LGF/high\_contrast/report\_hc\_summary/make\_fig.py}；附录表格由同目录的 \texttt{gen\_appendix.py} 直接从结果文件生成。
\end{itemize}

\clearpage
\newgeometry{margin=1.5cm}
\appendix
\section*{附录：全部实验数据}
\addcontentsline{toc}{section}{附录：全部实验数据}
下面按方法列出本文用到的全部实验数据，方法名统一用现在的叫法。表格由同目录的 \texttt{gen\_appendix.py} 从原报告的表格和结果文件生成，数值没有改动。不同附录之间判据不同，见各节开头。
\section{FD-LGF，教授的 $\chi^x$ 系列（二维、三维）}
数据来自 11 月 3 日的 $\chi^x$ 实验（\texttt{FD\_LGF/report/}）。算例 A：$a=\chi^x$；B：阶跃分层；C：$\chi^x$ 加规则夹杂；D：$\chi^x$ 加随机夹杂。$w$ 是局部精确解外环的宽度。A、B 没有特征区，只用背景逆，$w$ 不起作用：二维的结果填在 $w=1/16$ 一列，三维填在不带 $w$ 的“FD-LGF”一列。判据 PCG，$\|r\|/\|f\|\le10^{-10}$，单线程，3 次取中位数。
\subsection{二维}
\begin{table}[H]\centering\small
\caption{二维：PCG 迭代次数（$\chi\in\{10^2,10^4,10^6,10^8\}$ 的最小--最大值；$^{\times k}$ 表示有 $k$ 个 $\chi$ 在 1000 次内未收敛）}\label{tab:its2d}
\begin{adjustbox}{max width=\textwidth}
\begin{tabular}{llcccccc}
\toprule
算例 & $n$ & 传统 LGF & GMG & BoomerAMG & FD-LGF $w{=}\frac1{16}$ & FD-LGF $w{=}\frac1{32}$ & FD-LGF $w{=}\frac1{64}$ \\
\midrule
A: $\chi^x$ & 127 & 110--1000$^{\times2}$ & 7 & 7 & 2--3 &  &  \\
 & 255 & 117--1000$^{\times2}$ & 7--8 & 7 & 2--3 &  &  \\
 & 511 & 123--1000$^{\times2}$ & 7--8 & 7 & 2--3 &  &  \\
 & 1023 & 127--1000$^{\times3}$ & 7--8 & 7--8 & 2--3 &  &  \\
\midrule
B: 阶跃层 & 127 & 15--51 & 7 & 8 & 1--3 &  &  \\
 & 255 & 15--53 & 7 & 8 & 1--3 &  &  \\
 & 511 & 17--55 & 7 & 8 & 1--3 &  &  \\
 & 1023 & 19--58 & 7 & 9 & 1--3 &  &  \\
\midrule
C: $\chi^x$+规则夹杂 & 127 & 1000$^{\times4}$ & 48--58 & 9--10 & 8 & 10--11 & 14--15 \\
 & 255 & 1000$^{\times4}$ & 56--68 & 9--10 & 8 & 10--11 & 14--15 \\
 & 511 & 1000$^{\times4}$ & 82--87 & 10 & 8 & 10--11 & 13--14 \\
 & 1023 & 1000$^{\times4}$ & 90--95 & 10--11 & 7--8 & 10--11 & 13--14 \\
\midrule
D: $\chi^x$+随机夹杂 & 127 & 1000$^{\times4}$ & 73--80 & 8--9 & 9--10 & 13--15 & 19--21 \\
 & 255 & 1000$^{\times4}$ & 87--93 & 10 & 9--10 & 13--14 & 19--21 \\
 & 511 & 1000$^{\times4}$ & 98--104 & 10 & 9--10 & 13--14 & 19--20 \\
 & 1023 & 1000$^{\times4}$ & 108--114 & 10--11 & 8--10 & 13--14 & 19--20 \\
\bottomrule
\end{tabular}
\end{adjustbox}
\end{table}

\begin{table}[H]\centering\small
\caption{二维：误差 $\|u_k-u_h\|/\|u_h\|\le10^{-8}$ 所需迭代数（各 $\chi$ 的范围；--：1000 次内未达到）}\label{tab:k82d}
\begin{adjustbox}{max width=\textwidth}
\begin{tabular}{llcccccc}
\toprule
算例 & $n$ & 传统 LGF & GMG & BoomerAMG & FD-LGF $w{=}\frac1{16}$ & FD-LGF $w{=}\frac1{32}$ & FD-LGF $w{=}\frac1{64}$ \\
\midrule
A: $\chi^x$ & 127 & 68 & 5--7 & 5 & 2--3 &  &  \\
 & 255 & 71 & 5--7 & 5 & 2--3 &  &  \\
 & 511 & 73 & 5--6 & 5 & 2--3 &  &  \\
 & 1023 & 73 & 5--6 & 5 & 2--3 &  &  \\
\midrule
B: 阶跃层 & 127 & 11--69 & 5 & 6 & 3 &  &  \\
 & 255 & 11--71 & 5 & 5--6 & 3 &  &  \\
 & 511 & 10--72 & 5 & 5--6 & 3 &  &  \\
 & 1023 & 10--72 & 4--5 & 5--6 & 3 &  &  \\
\midrule
C: $\chi^x$+规则夹杂 & 127 & -- & 55--56 & 8--9 & 9--10 & 12--13 & 15--18 \\
 & 255 & -- & 64--65 & 8--9 & 9--10 & 12--13 & 15--17 \\
 & 511 & -- & 81--84 & 9 & 9--10 & 12--13 & 15--16 \\
 & 1023 & -- & 88--89 & 9 & 9--10 & 12--13 & 15--16 \\
\midrule
D: $\chi^x$+随机夹杂 & 127 & -- & 75--77 & 7--8 & 10--11 & 16--18 & 23--26 \\
 & 255 & -- & 88--89 & 8--9 & 11 & 16--18 & 22--26 \\
 & 511 & -- & 98--100 & 9--11 & 11 & 15--18 & 22--25 \\
 & 1023 & -- & 107--108 & 9 & 11 & 15--18 & 22--25 \\
\bottomrule
\end{tabular}
\end{adjustbox}
\end{table}

\begin{table}[H]\centering\small
\caption{二维：按残差规则 $\|r\|/\|f\|\le10^{-10}$ 停止时的代数误差 $\|u-u_h\|/\|u_h\|$（各 $\chi$ 中的最大值）}\label{tab:err2d}
\begin{adjustbox}{max width=\textwidth}
\begin{tabular}{llcccccc}
\toprule
算例 & $n$ & 传统 LGF & GMG & BoomerAMG & FD-LGF $w{=}\frac1{16}$ & FD-LGF $w{=}\frac1{32}$ & FD-LGF $w{=}\frac1{64}$ \\
\midrule
A: $\chi^x$ & 127 & $2.8{\times}10^{-2}$ & $3.2{\times}10^{-10}$ & $3.5{\times}10^{-12}$ & $1.4{\times}10^{-9}$ &  &  \\
 & 255 & $3.7{\times}10^{-2}$ & $3.0{\times}10^{-12}$ & $5.1{\times}10^{-12}$ & $1.4{\times}10^{-9}$ &  &  \\
 & 511 & $4.2{\times}10^{-2}$ & $1.9{\times}10^{-12}$ & $5.9{\times}10^{-12}$ & $1.4{\times}10^{-9}$ &  &  \\
 & 1023 & $4.5{\times}10^{-2}$ & $1.3{\times}10^{-12}$ & $1.1{\times}10^{-11}$ & $1.4{\times}10^{-9}$ &  &  \\
\midrule
B: 阶跃层 & 127 & $6.0{\times}10^{-7}$ & $3.3{\times}10^{-12}$ & $6.5{\times}10^{-12}$ & $3.7{\times}10^{-4}$ &  &  \\
 & 255 & $6.0{\times}10^{-7}$ & $3.0{\times}10^{-12}$ & $6.7{\times}10^{-13}$ & $3.8{\times}10^{-4}$ &  &  \\
 & 511 & $6.0{\times}10^{-7}$ & $2.6{\times}10^{-12}$ & $2.9{\times}10^{-12}$ & $3.8{\times}10^{-4}$ &  &  \\
 & 1023 & $6.0{\times}10^{-7}$ & $2.4{\times}10^{-12}$ & $1.1{\times}10^{-13}$ & $3.8{\times}10^{-4}$ &  &  \\
\midrule
C: $\chi^x$+规则夹杂 & 127 & $1.9{\times}10^{-1}$ & $1.2{\times}10^{-6}$ & $2.6{\times}10^{-9}$ & $2.5{\times}10^{-7}$ & $6.2{\times}10^{-7}$ & $4.0{\times}10^{-7}$ \\
 & 255 & $2.2{\times}10^{-1}$ & $2.7{\times}10^{-7}$ & $1.8{\times}10^{-9}$ & $2.7{\times}10^{-7}$ & $5.1{\times}10^{-7}$ & $2.2{\times}10^{-7}$ \\
 & 511 & $2.3{\times}10^{-1}$ & $9.9{\times}10^{-9}$ & $2.6{\times}10^{-10}$ & $2.3{\times}10^{-7}$ & $4.1{\times}10^{-7}$ & $7.9{\times}10^{-7}$ \\
 & 1023 & $2.3{\times}10^{-1}$ & $2.9{\times}10^{-9}$ & $9.2{\times}10^{-10}$ & $7.1{\times}10^{-6}$ & $3.8{\times}10^{-7}$ & $6.9{\times}10^{-7}$ \\
\midrule
D: $\chi^x$+随机夹杂 & 127 & $3.4{\times}10^{-1}$ & $3.6{\times}10^{-8}$ & $7.3{\times}10^{-10}$ & $7.6{\times}10^{-7}$ & $2.4{\times}10^{-6}$ & $1.8{\times}10^{-6}$ \\
 & 255 & $3.7{\times}10^{-1}$ & $1.9{\times}10^{-8}$ & $1.5{\times}10^{-9}$ & $7.7{\times}10^{-7}$ & $1.7{\times}10^{-6}$ & $1.5{\times}10^{-6}$ \\
 & 511 & $3.9{\times}10^{-1}$ & $8.2{\times}10^{-9}$ & $1.5{\times}10^{-8}$ & $6.0{\times}10^{-7}$ & $2.0{\times}10^{-6}$ & $1.3{\times}10^{-6}$ \\
 & 1023 & $4.0{\times}10^{-1}$ & $4.9{\times}10^{-9}$ & $2.5{\times}10^{-10}$ & $3.0{\times}10^{-6}$ & $1.9{\times}10^{-6}$ & $1.2{\times}10^{-6}$ \\
\bottomrule
\end{tabular}
\end{adjustbox}
\end{table}

\begin{table}[H]\centering\small
\caption{二维：总时间 setup+solve（$\chi=10^8$，单线程，ms）}\label{tab:time2d}
\begin{adjustbox}{max width=\textwidth}
\begin{tabular}{llcccccc}
\toprule
算例 & $n$ & 传统 LGF & GMG & BoomerAMG & FD-LGF $w{=}\frac1{16}$ & FD-LGF $w{=}\frac1{32}$ & FD-LGF $w{=}\frac1{64}$ \\
\midrule
A: $\chi^x$ & 127 & 542 & 16 & 18 & 4 &  &  \\
 & 255 & 2.36\,s & 55 & 77 & 15 &  &  \\
 & 511 & 11.07\,s & 219 & 306 & 62 &  &  \\
 & 1023 & 46.30\,s & 899 & 1.14\,s & 251 &  &  \\
\midrule
B: 阶跃层 & 127 & 26 & 15 & 21 & 3 &  &  \\
 & 255 & 114 & 49 & 83 & 12 &  &  \\
 & 511 & 536 & 198 & 323 & 50 &  &  \\
 & 1023 & 3.04\,s & 941 & 1.35\,s & 195 &  &  \\
\midrule
C: $\chi^x$+规则夹杂 & 127 & 532 & 61 & 27 & 31 & 27 & 29 \\
 & 255 & 2.11\,s & 250 & 103 & 150 & 108 & 110 \\
 & 511 & 8.09\,s & 1.37\,s & 376 & 806 & 482 & 416 \\
 & 1023 & 38.69\,s & 8.20\,s & 1.53\,s & 4.70\,s & 2.59\,s & 2.26\,s \\
\midrule
D: $\chi^x$+随机夹杂 & 127 & 627 & 87 & 27 & 36 & 35 & 40 \\
 & 255 & 2.22\,s & 391 & 109 & 160 & 135 & 143 \\
 & 511 & 8.01\,s & 1.80\,s & 444 & 891 & 640 & 597 \\
 & 1023 & 37.92\,s & 9.33\,s & 1.62\,s & 5.02\,s & 3.15\,s & 2.72\,s \\
\bottomrule
\end{tabular}
\end{adjustbox}
\end{table}

\begin{table}[H]\centering\small
\caption{二维：误差降到 $10^{-8}$ 的时间 setup$+k_8\times$单步时间（$\chi=10^8$，ms）}\label{tab:t82d}
\begin{adjustbox}{max width=\textwidth}
\begin{tabular}{llcccccc}
\toprule
算例 & $n$ & 传统 LGF & GMG & BoomerAMG & FD-LGF $w{=}\frac1{16}$ & FD-LGF $w{=}\frac1{32}$ & FD-LGF $w{=}\frac1{64}$ \\
\midrule
A: $\chi^x$ & 127 & -- & 16 & 15 & 4 &  &  \\
 & 255 & -- & 50 & 65 & 15 &  &  \\
 & 511 & -- & 184 & 260 & 62 &  &  \\
 & 1023 & -- & 752 & 974 & 251 &  &  \\
\midrule
B: 阶跃层 & 127 & 40 & 12 & 18 & 7 &  &  \\
 & 255 & 171 & 41 & 66 & 25 &  &  \\
 & 511 & 785 & 163 & 258 & 102 &  &  \\
 & 1023 & 4.20\,s & 766 & 1.01\,s & 409 &  &  \\
\midrule
C: $\chi^x$+规则夹杂 & 127 & -- & 70 & 27 & 35 & 32 & 34 \\
 & 255 & -- & 283 & 103 & 168 & 127 & 126 \\
 & 511 & -- & 1.37\,s & 355 & 882 & 558 & 476 \\
 & 1023 & -- & 8.12\,s & 1.44\,s & 5.24\,s & 2.97\,s & 2.56\,s \\
\midrule
D: $\chi^x$+随机夹杂 & 127 & -- & 89 & 25 & 41 & 43 & 51 \\
 & 255 & -- & 395 & 103 & 178 & 168 & 183 \\
 & 511 & -- & 1.80\,s & 469 & 982 & 791 & 733 \\
 & 1023 & -- & 9.24\,s & 1.53\,s & 5.64\,s & 3.82\,s & 3.31\,s \\
\bottomrule
\end{tabular}
\end{adjustbox}
\end{table}

\begin{table}[H]\centering\small
\caption{二维，算例 A（$n=1023$）：相对连续解的总误差 $\|u-u^*\|/\|u^*\|$}\label{tab:aerr二维}
\begin{tabular}{lcccc}
\toprule
$\chi$ & 传统 LGF & GMG & BoomerAMG & FD-LGF $w{=}\frac1{16}$ \\
\midrule
$10^{2}$ & $1.7{\times}10^{-6}$ & $1.7{\times}10^{-6}$ & $1.7{\times}10^{-6}$ & $1.7{\times}10^{-6}$ \\
$10^{4}$ & $1.0{\times}10^{-5}$ & $1.0{\times}10^{-5}$ & $1.0{\times}10^{-5}$ & $1.0{\times}10^{-5}$ \\
$10^{6}$ & $2.1{\times}10^{-3}$ & $2.5{\times}10^{-5}$ & $2.5{\times}10^{-5}$ & $2.5{\times}10^{-5}$ \\
$10^{8}$ & $4.5{\times}10^{-2}$ & $4.7{\times}10^{-5}$ & $4.7{\times}10^{-5}$ & $4.7{\times}10^{-5}$ \\
\bottomrule
\end{tabular}
\end{table}
\subsection{三维}
\begin{table}[H]\centering\small
\caption{三维：PCG 迭代次数（$\chi\in\{10^2,10^4,10^6,10^8\}$ 的最小--最大值；$^{\times k}$ 表示有 $k$ 个 $\chi$ 在 500 次内未收敛）}\label{tab:its3d}
\begin{adjustbox}{max width=\textwidth}
\begin{tabular}{llcccccc}
\toprule
算例 & $n$ & 传统 LGF & GMG & BoomerAMG & FD-LGF & FD-LGF $w{=}\frac1{32}$ & FD-LGF $w{=}\frac1{64}$ \\
\midrule
A: $\chi^x$ & 31 & 68--500$^{\times2}$ & 8--9 & 5--7 & 2--3 &  &  \\
 & 63 & 97--500$^{\times3}$ & 9 & 7 & 2--3 &  &  \\
 & 127 & 111--500$^{\times3}$ & 9 & 7--9 & 2--3 &  &  \\
\midrule
B: 阶跃层 & 31 & 20--48 & 9 & 7--8 & 1--3 &  &  \\
 & 63 & 19--59 & 9 & 9 & 1--3 &  &  \\
 & 127 & 19--62 & 9 & 10--11 & 1--3 &  &  \\
\midrule
C: $\chi^x$+规则夹杂 & 31 & 500$^{\times4}$ & 57--62 & 11 &  & 12 & 12 \\
 & 63 & 500$^{\times4}$ & 75--89 & 12--14 &  & 11 & 15--16 \\
 & 127 & 500$^{\times4}$ & 89--103 & 14--15 &  & 10 & 13--14 \\
\midrule
D: $\chi^x$+随机夹杂 & 31 & 500$^{\times4}$ & 70--88 & 13--14 &  & 12--13 & 12--13 \\
 & 63 & 500$^{\times4}$ & 100--107 & 15--16 &  & 12 & 17 \\
 & 127 & 500$^{\times4}$ & 123--135 & 15--17 &  & 11 & 16 \\
\bottomrule
\end{tabular}
\end{adjustbox}
\end{table}

\begin{table}[H]\centering\small
\caption{三维：误差 $\|u_k-u_h\|/\|u_h\|\le10^{-8}$ 所需迭代数（各 $\chi$ 的范围；--：300 次内未达到）}\label{tab:k83d}
\begin{adjustbox}{max width=\textwidth}
\begin{tabular}{llcccccc}
\toprule
算例 & $n$ & 传统 LGF & GMG & BoomerAMG & FD-LGF & FD-LGF $w{=}\frac1{32}$ & FD-LGF $w{=}\frac1{64}$ \\
\midrule
A: $\chi^x$ & 31 & 54--244 & 6--8 & 5 & 2--3 &  &  \\
 & 63 & 65 & 6--7 & 5--6 & 2--3 &  &  \\
 & 127 & 70 & 6--7 & 6--8 & 2--3 &  &  \\
\midrule
B: 阶跃层 & 31 & 12--80 & 6 & 6 & 3 &  &  \\
 & 63 & 12--85 & 6 & 7--8 & 3 &  &  \\
 & 127 & 12--87 & 5 & 8--9 & 3 &  &  \\
\midrule
C: $\chi^x$+规则夹杂 & 31 & -- & 64--67 & 10--11 &  & 13--14 & 13--14 \\
 & 63 & -- & 84--86 & 11--14 &  & 12--14 & 15--18 \\
 & 127 & -- & 99--100 & 13--14 &  & 10--13 & 14--16 \\
\midrule
D: $\chi^x$+随机夹杂 & 31 & -- & 80--86 & 12--14 &  & 13--15 & 13--15 \\
 & 63 & -- & 104--110 & 13--16 &  & 12--13 & 17--19 \\
 & 127 & -- & 123--129 & 15 &  & 11--12 & 16--18 \\
\bottomrule
\end{tabular}
\end{adjustbox}
\end{table}

\begin{table}[H]\centering\small
\caption{三维：按残差规则 $\|r\|/\|f\|\le10^{-10}$ 停止时的代数误差 $\|u-u_h\|/\|u_h\|$（各 $\chi$ 中的最大值）}\label{tab:err3d}
\begin{adjustbox}{max width=\textwidth}
\begin{tabular}{llcccccc}
\toprule
算例 & $n$ & 传统 LGF & GMG & BoomerAMG & FD-LGF & FD-LGF $w{=}\frac1{32}$ & FD-LGF $w{=}\frac1{64}$ \\
\midrule
A: $\chi^x$ & 31 & $2.4{\times}10^{-2}$ & $6.1{\times}10^{-9}$ & $1.7{\times}10^{-10}$ & $3.3{\times}10^{-9}$ &  &  \\
 & 63 & $5.8{\times}10^{-2}$ & $7.1{\times}10^{-11}$ & $1.5{\times}10^{-10}$ & $3.3{\times}10^{-9}$ &  &  \\
 & 127 & $8.4{\times}10^{-2}$ & $3.2{\times}10^{-11}$ & $1.5{\times}10^{-10}$ & $3.3{\times}10^{-9}$ &  &  \\
\midrule
B: 阶跃层 & 31 & $7.3{\times}10^{-7}$ & $3.3{\times}10^{-12}$ & $4.1{\times}10^{-10}$ & $2.0{\times}10^{-4}$ &  &  \\
 & 63 & $7.5{\times}10^{-7}$ & $8.9{\times}10^{-13}$ & $1.6{\times}10^{-10}$ & $2.0{\times}10^{-4}$ &  &  \\
 & 127 & $7.5{\times}10^{-7}$ & $3.4{\times}10^{-13}$ & $3.9{\times}10^{-10}$ & $2.0{\times}10^{-4}$ &  &  \\
\midrule
C: $\chi^x$+规则夹杂 & 31 & $3.3{\times}10^{-1}$ & $1.0{\times}10^{-7}$ & $3.3{\times}10^{-9}$ &  & $2.7{\times}10^{-7}$ & $2.7{\times}10^{-7}$ \\
 & 63 & $4.1{\times}10^{-1}$ & $2.2{\times}10^{-7}$ & $1.5{\times}10^{-8}$ &  & $4.5{\times}10^{-7}$ & $3.1{\times}10^{-7}$ \\
 & 127 & $4.4{\times}10^{-1}$ & $2.9{\times}10^{-7}$ & $2.4{\times}10^{-9}$ &  & $2.5{\times}10^{-7}$ & $1.2{\times}10^{-7}$ \\
\midrule
D: $\chi^x$+随机夹杂 & 31 & $5.4{\times}10^{-1}$ & $3.0{\times}10^{-7}$ & $6.6{\times}10^{-9}$ &  & $7.0{\times}10^{-8}$ & $7.0{\times}10^{-8}$ \\
 & 63 & $6.3{\times}10^{-1}$ & $6.0{\times}10^{-8}$ & $1.3{\times}10^{-8}$ &  & $1.2{\times}10^{-8}$ & $7.9{\times}10^{-8}$ \\
 & 127 & $6.6{\times}10^{-1}$ & $9.9{\times}10^{-9}$ & $4.4{\times}10^{-9}$ &  & $3.2{\times}10^{-8}$ & $1.7{\times}10^{-7}$ \\
\bottomrule
\end{tabular}
\end{adjustbox}
\end{table}

\begin{table}[H]\centering\small
\caption{三维：总时间 setup+solve（$\chi=10^8$，单线程，ms）}\label{tab:time3d}
\begin{adjustbox}{max width=\textwidth}
\begin{tabular}{llcccccc}
\toprule
算例 & $n$ & 传统 LGF & GMG & BoomerAMG & FD-LGF & FD-LGF $w{=}\frac1{32}$ & FD-LGF $w{=}\frac1{64}$ \\
\midrule
A: $\chi^x$ & 31 & 745 & 33 & 100 & 5 &  &  \\
 & 63 & 6.52\,s & 257 & 1.23\,s & 41 &  &  \\
 & 127 & 47.96\,s & 2.62\,s & 8.86\,s & 335 &  &  \\
\midrule
B: 阶跃层 & 31 & 69 & 31 & 119 & 4 &  &  \\
 & 63 & 698 & 265 & 1.21\,s & 29 &  &  \\
 & 127 & 5.14\,s & 2.28\,s & 12.80\,s & 229 &  &  \\
\midrule
C: $\chi^x$+规则夹杂 & 31 & 688 & 137 & 134 &  & 48 & 67 \\
 & 63 & 6.25\,s & 1.49\,s & 1.37\,s &  & 503 & 429 \\
 & 127 & 44.08\,s & 13.50\,s & 10.09\,s &  & 21.21\,s & 9.25\,s \\
\midrule
D: $\chi^x$+随机夹杂 & 31 & 710 & 164 & 155 &  & 41 & 42 \\
 & 63 & 6.29\,s & 1.93\,s & 1.66\,s &  & 369 & 376 \\
 & 127 & 43.00\,s & 19.73\,s & 11.66\,s &  & 9.80\,s & 6.01\,s \\
\bottomrule
\end{tabular}
\end{adjustbox}
\end{table}

\begin{table}[H]\centering\small
\caption{三维：误差降到 $10^{-8}$ 的时间 setup$+k_8\times$单步时间（$\chi=10^8$，ms）}\label{tab:t83d}
\begin{adjustbox}{max width=\textwidth}
\begin{tabular}{llcccccc}
\toprule
算例 & $n$ & 传统 LGF & GMG & BoomerAMG & FD-LGF & FD-LGF $w{=}\frac1{32}$ & FD-LGF $w{=}\frac1{64}$ \\
\midrule
A: $\chi^x$ & 31 & -- & 33 & 100 & 5 &  &  \\
 & 63 & -- & 219 & 1.15\,s & 41 &  &  \\
 & 127 & -- & 2.25\,s & 8.40\,s & 335 &  &  \\
\midrule
B: 阶跃层 & 31 & 115 & 24 & 112 & 10 &  &  \\
 & 63 & 1.12\,s & 205 & 1.15\,s & 78 &  &  \\
 & 127 & 7.98\,s & 1.61\,s & 11.56\,s & 611 &  &  \\
\midrule
C: $\chi^x$+规则夹杂 & 31 & -- & 153 & 127 &  & 54 & 75 \\
 & 63 & -- & 1.66\,s & 1.43\,s &  & 576 & 489 \\
 & 127 & -- & 14.95\,s & 9.66\,s &  & 22.29\,s & 9.99\,s \\
\midrule
D: $\chi^x$+随机夹杂 & 31 & -- & 188 & 155 &  & 49 & 50 \\
 & 63 & -- & 2.07\,s & 1.73\,s &  & 390 & 412 \\
 & 127 & -- & 19.73\,s & 11.66\,s &  & 10.07\,s & 6.39\,s \\
\bottomrule
\end{tabular}
\end{adjustbox}
\end{table}

\begin{table}[H]\centering\small
\caption{三维，算例 A（$n=127$）：相对连续解的总误差 $\|u-u^*\|/\|u^*\|$}\label{tab:aerr三维}
\begin{tabular}{lcccc}
\toprule
$\chi$ & 传统 LGF & GMG & BoomerAMG & FD-LGF \\
\midrule
$10^{2}$ & $6.3{\times}10^{-5}$ & $6.3{\times}10^{-5}$ & $6.3{\times}10^{-5}$ & $6.3{\times}10^{-5}$ \\
$10^{4}$ & $4.8{\times}10^{-4}$ & $4.8{\times}10^{-4}$ & $4.8{\times}10^{-4}$ & $4.8{\times}10^{-4}$ \\
$10^{6}$ & $9.3{\times}10^{-3}$ & $1.3{\times}10^{-3}$ & $1.3{\times}10^{-3}$ & $1.3{\times}10^{-3}$ \\
$10^{8}$ & $8.5{\times}10^{-2}$ & $2.5{\times}10^{-3}$ & $2.5{\times}10^{-3}$ & $2.5{\times}10^{-3}$ \\
\bottomrule
\end{tabular}
\end{table}
\FloatBarrier
\section{FD-LGF，更一般的 10 类二维和 5 类三维系数}
数据来自报告 01（\texttt{FD\_LGF/high\_contrast/report\_hc/}）。所有系数写成 $a=\chi^g$，$\chi\in\{10^2,10^4,10^6,10^8\}$。二维 $n=1023$ 为确认实验，方案锁定 \texttt{c2ff7672}；其余为探索性扫描。
\begin{table}[H]\centering\small
\caption{探索性扫描（$n=255, 511$）：迭代次数范围}\label{tab:scan}
\begin{tabular}{lcccc}
\toprule
系数 & BoomerAMG & GMG & FD-LGF $w{=}\frac1{16}$ & FD-LGF $w{=}\frac1{64}$ \\
\midrule
sep2 可分离 & 7--9 & 8 & 3--4 & 3--4 \\
chan 通道 & 10--12 & 19--32 & 9--11 & 17--25 \\
checker 棋盘 & 9--10 & 19--24 & 3--10 & 6--21 \\
lognorm 对数正态 & 7--9 & 9--17 & 5--22 & 24--28 \\
layered 类SPE10分层 & 9--12 & 13--32 & 16--22 & 22--24 \\
diag 斜条纹 & 10 & 29--38 & 1 & 4--26 \\
walls 绝缘墙 & 9--10 & 16--20 & 9 & 14 \\
perc 随机两相 & 11--12 & 37--101 & 1 & 11--26 \\
diagfine 细斜条纹 & 8--13 & 43--398 & 1 & 1 \\
rough 粗糙随机场 & 8--12 & 10--54 & 1--17 & 12--22 \\
\bottomrule
\end{tabular}
\end{table}

\begin{table}[H]\centering\small
\caption{三维探索性扫描：迭代次数（$n=63$）}\label{tab:s3its63}
\begin{tabular}{lccccc}
\toprule
系数 & BoomerAMG $\theta{=}0.25$ & BoomerAMG $\theta{=}0.5$ & GMG & FD-LGF $w{=}\frac1{32}$ & FD-LGF $w{=}\frac1{64}$ \\
\midrule
sheets 薄层 & 9 & 8--9 & 16--18 & 2--3 & 2--3 \\
tubes 管道网 & 9 & 10 & 15--20 & 5--13 & 7--18 \\
insul 绝缘板 & 8 & 9 & 9 & 7 & 8 \\
perc3 随机两相 & 8 & 10 & 19--25 & -- & -- \\
lognorm3 对数正态 & 7--8 & 10--11 & 10--16 & -- & -- \\
\bottomrule
\end{tabular}
\end{table}

\begin{table}[H]\centering\small
\caption{三维探索性扫描：迭代次数（$n=127$）}\label{tab:s3its127}
\begin{tabular}{lccccc}
\toprule
系数 & BoomerAMG $\theta{=}0.25$ & BoomerAMG $\theta{=}0.5$ & GMG & FD-LGF $w{=}\frac1{32}$ & FD-LGF $w{=}\frac1{64}$ \\
\midrule
sheets 薄层 & 9 & 10--11 & 14--16 & 2--3 & 2--3 \\
tubes 管道网 & 10 & 12 & 19--26 & 5--12 & 7--17 \\
insul 绝缘板 & 8 & 11--12 & 10 & 8 & 9 \\
perc3 随机两相 & 9 & 12 & 21--28 & -- & -- \\
lognorm3 对数正态 & 8 & 11--12 & 10--11 & -- & -- \\
\bottomrule
\end{tabular}
\end{table}

\begin{table}[H]\centering\small
\caption{三维探索性扫描：总时间（s，$n=127$，约 $2\times10^6$ 未知量）}\label{tab:s3time}
\begin{tabular}{lccccc}
\toprule
系数 & BoomerAMG $\theta{=}0.25$ & BoomerAMG $\theta{=}0.5$ & GMG & FD-LGF $w{=}\frac1{32}$ & FD-LGF $w{=}\frac1{64}$ \\
\midrule
sheets 薄层 & 12.2--14.0 & 12.2--14.3 & 3.1--3.4 & 0.3--0.5 & 0.3--0.5 \\
tubes 管道网 & 14.1--16.8 & 12.8--16.2 & 4.1--5.2 & 22.6--27.1 & 4.6--8.5 \\
insul 绝缘板 & 7.9--9.3 & 12.1--12.5 & 2.1--2.4 & 733.1--791.8 & 115.4--115.9 \\
perc3 随机两相 & 12.6--12.7 & 11.4--12.7 & 3.6--4.6 & -- & -- \\
lognorm3 对数正态 & 8.0--8.4 & 10.5--12.8 & 2.0--2.1 & -- & -- \\
\bottomrule
\end{tabular}
\end{table}
\begin{table}[H]\centering\small
\caption{$n=1023$：PCG 迭代次数（4 个 $\chi$ 的范围）}\label{tab:its}
\begin{tabular}{lcccc}
\toprule
系数 & BoomerAMG & GMG & FD-LGF $w{=}\frac1{16}$ & FD-LGF $w{=}\frac1{64}$ \\
\midrule
sep2 可分离 & 8--9 & 8 & 3--4 & 3--4 \\
chan 通道 & 12--13 & 21--34 & 9--11 & 17--25 \\
checker 棋盘 & 11 & 23--27 & 3--10 & 6--21 \\
lognorm 对数正态 & 8--9 & 10--17 & 6--22 & 25--28 \\
layered 类SPE10分层 & 9--12 & 14--33 & 16--22 & 23--24 \\
diag 斜条纹 & 10--11 & 30--45 & 1 & 4--26 \\
walls 绝缘墙 & 10--11 & 19--21 & 9 & 14 \\
perc 随机两相 & 12--13 & 37--106 & 1 & 11--26 \\
diagfine 细斜条纹 & 13 & 47--368 & 1 & 1 \\
rough 粗糙随机场 & 9--11 & 11--54 & 1--18 & 16--23 \\
\bottomrule
\end{tabular}
\end{table}

\begin{table}[H]\centering\small
\caption{$n=1023$：误差降到 $10^{-8}$ 所需迭代数}\label{tab:k8}
\begin{tabular}{lcccc}
\toprule
系数 & BoomerAMG & GMG & FD-LGF $w{=}\frac1{16}$ & FD-LGF $w{=}\frac1{64}$ \\
\midrule
sep2 可分离 & 5--6 & 4--6 & 3--4 & 3--4 \\
chan 通道 & 9--10 & 15--30 & 9--11 & 14--25 \\
checker 棋盘 & 8 & 15--23 & 8 & 17 \\
lognorm 对数正态 & 5--7 & 6--16 & 5--13 & 16--22 \\
layered 类SPE10分层 & 6--9 & 9--27 & 12--14 & 14--18 \\
diag 斜条纹 & 8 & 24--43 & 1 & 18--20 \\
walls 绝缘墙 & 7--8 & 13--16 & 7 & 11 \\
perc 随机两相 & 9--10 & 27--99 & 1--2 & 20--26 \\
diagfine 细斜条纹 & 9--10 & 34--282 & 1 & 1 \\
rough 粗糙随机场 & 6--9 & 7--44 & 1--10 & 12--14 \\
\bottomrule
\end{tabular}
\end{table}

\begin{table}[H]\centering\small
\caption{$n=1023$：总时间 setup+solve（s，单线程）}\label{tab:time}
\begin{tabular}{lcccc}
\toprule
系数 & BoomerAMG & GMG & FD-LGF $w{=}\frac1{16}$ & FD-LGF $w{=}\frac1{64}$ \\
\midrule
sep2 可分离 & 1.18--1.28 & 0.78--0.87 & 1.17--1.38 & 1.16--1.38 \\
chan 通道 & 1.54--1.68 & 1.63--2.34 & 5.01--5.50 & 2.40--3.21 \\
checker 棋盘 & 1.42--1.47 & 1.79--2.03 & 6.59--8.25 & 4.28--6.61 \\
lognorm 对数正态 & 1.18--1.26 & 1.01--1.32 & 10.96--22.88 & 8.53--17.12 \\
layered 类SPE10分层 & 1.26--1.50 & 1.24--2.18 & 8.63--20.35 & 6.55--14.08 \\
diag 斜条纹 & 1.43--1.51 & 2.22--3.06 & 14.37--14.57 & 2.62--6.07 \\
walls 绝缘墙 & 1.41--1.51 & 1.40--1.47 & 4.74--4.80 & 1.60--1.71 \\
perc 随机两相 & 1.72--1.82 & 2.46--7.02 & 14.35--14.43 & 9.77--13.22 \\
diagfine 细斜条纹 & 1.85--1.89 & 3.31--23.28 & 14.31--14.44 & 15.23--15.51 \\
rough 粗糙随机场 & 1.23--1.50 & 1.04--3.74 & 14.99--24.53 & 8.83--28.66 \\
\bottomrule
\end{tabular}
\end{table}

\begin{table}[H]\centering\small
\caption{$n=1023$：setup 时间（s）}\label{tab:setup}
\begin{tabular}{lcccc}
\toprule
系数 & BoomerAMG & GMG & FD-LGF $w{=}\frac1{16}$ & FD-LGF $w{=}\frac1{64}$ \\
\midrule
sep2 可分离 & 0.55 & 0.28--0.30 & 0.52--0.53 & 0.51--0.52 \\
chan 通道 & 0.59--0.61 & 0.29--0.30 & 3.32--3.34 & 0.68--0.71 \\
checker 棋盘 & 0.56--0.57 & 0.29--0.30 & 5.76--5.86 & 2.91--3.09 \\
lognorm 对数正态 & 0.53--0.55 & 0.29--0.30 & 3.47--15.84 & 1.72--5.74 \\
layered 类SPE10分层 & 0.54 & 0.28--0.29 & 2.60--13.12 & 0.99--4.81 \\
diag 斜条纹 & 0.61--0.62 & 0.29--0.30 & 14.02--14.07 & 1.75--1.91 \\
walls 绝缘墙 & 0.60--0.61 & 0.28--0.29 & 3.10--3.13 & 0.52--0.54 \\
perc 随机两相 & 0.68--0.70 & 0.28--0.29 & 13.99--14.05 & 6.45--6.89 \\
diagfine 细斜条纹 & 0.68--0.69 & 0.29--0.32 & 13.96--14.08 & 14.58--14.73 \\
rough 粗糙随机场 & 0.54--0.57 & 0.29--0.30 & 14.49--16.14 & 1.91--19.39 \\
\bottomrule
\end{tabular}
\end{table}

\begin{table}[H]\centering\small
\caption{$n=1023$：FD-LGF 局部 LU 区域占网格的比例}\label{tab:loc}
\begin{tabular}{lcc}
\toprule
系数 & FD-LGF $w{=}\frac1{16}$ & FD-LGF $w{=}\frac1{64}$ \\
\midrule
sep2 可分离 & 0.00 & 0.00 \\
chan 通道 & 0.65 & 0.26 \\
checker 棋盘 & 0.80 & 0.59 \\
lognorm 对数正态 & 0.64--1.00 & 0.36--0.89 \\
layered 类SPE10分层 & 0.46--0.97 & 0.20--0.81 \\
diag 斜条纹 & 1.00 & 0.67 \\
walls 绝缘墙 & 0.64 & 0.24 \\
perc 随机两相 & 1.00 & 0.90 \\
diagfine 细斜条纹 & 1.00 & 1.00 \\
rough 粗糙随机场 & 1.00 & 0.65--1.00 \\
\bottomrule
\end{tabular}
\end{table}

\begin{table}[H]\centering\small
\caption{$n=1023$：按残差规则停止时的代数误差（各 $\chi$ 最大值）}\label{tab:err}
\begin{tabular}{lcccc}
\toprule
系数 & BoomerAMG & GMG & FD-LGF $w{=}\frac1{16}$ & FD-LGF $w{=}\frac1{64}$ \\
\midrule
sep2 可分离 & $7.3{\times}10^{-13}$ & $8.2{\times}10^{-14}$ & $2.2{\times}10^{-9}$ & $2.2{\times}10^{-9}$ \\
chan 通道 & $1.9{\times}10^{-11}$ & $2.2{\times}10^{-10}$ & $5.9{\times}10^{-7}$ & $1.7{\times}10^{-7}$ \\
checker 棋盘 & $3.6{\times}10^{-13}$ & $3.1{\times}10^{-9}$ & $1.4{\times}10^{-4}$ & $6.1{\times}10^{-5}$ \\
lognorm 对数正态 & $1.2{\times}10^{-12}$ & $8.8{\times}10^{-12}$ & $3.6{\times}10^{-11}$ & $2.0{\times}10^{-12}$ \\
layered 类SPE10分层 & $2.6{\times}10^{-12}$ & $1.1{\times}10^{-11}$ & $4.9{\times}10^{-12}$ & $5.4{\times}10^{-12}$ \\
diag 斜条纹 & $1.1{\times}10^{-11}$ & $1.4{\times}10^{-8}$ & $3.4{\times}10^{-12}$ & $6.7{\times}10^{-5}$ \\
walls 绝缘墙 & $2.8{\times}10^{-12}$ & $4.6{\times}10^{-10}$ & $8.2{\times}10^{-14}$ & $6.8{\times}10^{-13}$ \\
perc 随机两相 & $1.5{\times}10^{-11}$ & $1.6{\times}10^{-11}$ & $2.2{\times}10^{-7}$ & $6.4{\times}10^{-5}$ \\
diagfine 细斜条纹 & $1.1{\times}10^{-11}$ & $3.1{\times}10^{-10}$ & $5.5{\times}10^{-12}$ & $5.5{\times}10^{-12}$ \\
rough 粗糙随机场 & $1.9{\times}10^{-11}$ & $1.2{\times}10^{-10}$ & $4.9{\times}10^{-12}$ & $4.1{\times}10^{-11}$ \\
\bottomrule
\end{tabular}
\end{table}
\FloatBarrier
\section{bt-fd（三温）}
判据：双双精度原残差 $\|b-Ax\|/\|b\|<10^{-8}$，时间含构建，单线程，独立进程。
\subsection{H57 相同精度多方法锁定实验（sha256 \texttt{82dbdb84}）}
\begin{table}[H]
\centering\small
\begin{tabular}{lrrrrrr}
\toprule
真实 SC26\_01（5 轮） & 总时间中位数 & 范围 & setup & 求解 & 内层迭代 & /基线（配对中位数）\\
\midrule
amg-base & 33.00 s & 32.3--35.9 & 6.33 & 21.12 & 37 & 1 \\
amg-tuned & 31.86 & 30.5--32.6 & 6.65 & 19.19 & 34 & 0.925 \\
bt-v & 29.80 & 27.6--32.7 & 6.15 & 17.69 & 38 & 0.853 \\
bt-j & 35.21 & 33.9--40.0 & 6.81 & 22.53 & 38 & 1.067 \\
bt-fd LU，容差 $10^{-6}$ & 19.58 & 18.9--22.5 & 6.39 & 8.25 & 16 & 0.602 \\
bt-fd LU，容差 0.1 & 20.21 & 19.1--20.5 & 5.95 & 8.87 & 17 & -- \\
bt-fd Cholesky，容差 $10^{-6}$ & \textbf{16.78} & 16.4--18.9 & 4.99 & 7.18 & 16 & \textbf{0.508} \\
bt-fd Cholesky，容差 0.1 & 17.34 & 16.6--17.8 & 4.64 & 7.95 & 18 & -- \\
\bottomrule
\end{tabular}
\caption{全部 3 步修正（残差计算约 5--6\,s）。公共准备约 0.65\,s。除 bt-j 有一次写出 0.86\,GB swap 外，其余 swap 写出为 0。}
\end{table}
\begin{table}[H]
\centering\footnotesize\setlength{\tabcolsep}{4pt}
\begin{tabular}{lrrrrrrrrr}
\toprule
序列 & 特征区 & amg-base & amg-tuned & bt-v & bt-j & bt-fd & chol & chol/tuned & bt-v/tuned \\
\midrule
G1 球 & 4.7\% & 246.9 & 220.2 & 208.7 & 264.9 & 126.4 & \textbf{117.2} & 0.53 & 0.95 \\
G4 长椭球 & 9.1\% & 250.2 & 219.0 & 199.0 & 255.8 & 193.5 & \textbf{141.8} & 0.65 & 0.91 \\
G7 3 胶囊 & 6.6\% & 258.6 & 215.0 & 225.7 & 275.3 & 165.0 & \textbf{156.8} & 0.73 & 1.05 \\
G9 10 胶囊 & 16.3\% & 245.5 & 204.7 & \textbf{189.6} & 238.2 & 251.9 & 210.2 & 1.03 & 0.93 \\
G8 10 胶囊 & 18.4\% & 234.0 & 209.7 & \textbf{194.0} & 230.1 & 326.8 & 283.5 & 1.35 & 0.93 \\
G10 10 扁方块 & 20.3\% & 233.3 & 207.8 & \textbf{188.4} & 239.8 & 325.4 & 257.7 & 1.24 & 0.91 \\
\bottomrule
\end{tabular}
\caption{10 步总时间（s）。一遍，每步最多 10 次修正，实际 2--3 次。}
\end{table}
\subsection{审计修复后的 SC26\_01 配对复测和 H56 G8--G11 重跑（报告 24）}
\begin{table}[H]
\centering\footnotesize\setlength{\tabcolsep}{4pt}
\begin{tabular}{lrrrrrrrr}
\toprule
方法 & 总时间中位数 & 范围 & setup & 求解 & DD 残差 & 内层迭代 & 最终原残差 & 峰值内存 \\
\midrule
BoomerAMG 基线 ($\theta=0.5$) & 36.87\,s & 34.7--39.6 & 6.74 & 23.07 & 6.33 & 37 & $4.5\times10^{-10}$ & 4.0\,GB \\
BoomerAMG 筛选 ($\theta=0.25$) & 34.76\,s & 33.5--36.4 & 7.18 & 21.05 & 6.31 & 34 & $2.4\times10^{-10}$ & 4.1\,GB \\
bt-fd，MUMPS LU & 22.45\,s & 20.1--25.5 & 6.65 & 9.59 & 6.19 & 16 & $5.6\times10^{-11}$ & 4.3\,GB \\
bt-fd，CHOLMOD & 18.23\,s & 16.9--20.0 & 5.25 & 7.68 & 5.09 & 16 & $1.2\times10^{-10}$ & 4.8\,GB \\
\bottomrule
\end{tabular}
\caption{SC26\_01，$n=6\,291\,456$，20 次运行全部达标，内层 reason 全为 2，都是 3 步修正。公共准备（DD 表）另需约 0.7\,s。}
\end{table}
\begin{table}[H]
\centering\small
\begin{tabular}{lll}
\toprule
配对（同一轮内相除） & 中位数 & 5 轮范围 \\
\midrule
bt-fd LU / 基线 AMG & 0.572 & 0.562--0.734 \\
bt-fd LU / 筛选 AMG & 0.646 & 0.594--0.700 \\
bt-fd Cholesky / bt-fd LU & 0.812 & 0.737--0.908 \\
bt-fd Cholesky / 筛选 AMG & 0.525 & 0.500--0.548 \\
筛选 AMG / 基线 AMG & 0.947 & 0.878--1.048 \\
\bottomrule
\end{tabular}
\end{table}
\begin{table}[H]
\centering\scriptsize\setlength{\tabcolsep}{3pt}
\begin{tabular}{lrrrrrrrrr}
\toprule
序列 & 特征区 & 因子非零元（每块） & bt-fd 总时间 & AMG 总时间 & 时间比 & H56 原比值 & bt-fd 单步 & 峰值内存 & swap 写出 \\
\midrule
G8 & 18.4\% & $2.1\times10^{8}$ & 200.0\,s (330.1) & 119.5\,s (194.3) & \textbf{1.67} & 1.70 & 15--30\,s & 7.3\,GB & 0.4\,GB \\
G9 & 16.3\% & $1.4\times10^{8}$ & 130.3 (156.9) & 120.1 (139.3) & \textbf{1.09} & 1.13 & 11--13 & 6.2 & 0.4 \\
G10 & 20.3\% & $2.0\times10^{8}$ & 169.9 (349.3) & 113.5 (152.9) & \textbf{1.50} & 2.28 & 14--17 & 7.3 & 0 \\
G11 & 30.8\% & $3.6\times10^{8}$ & 783.7 (1209.6) & 116.7 (275.8) & \textbf{6.71} & 4.39 & 40--121 & 8.6 & 36\,GB \\
\bottomrule
\end{tabular}
\caption{括号内是 H56 原始值。峰值内存和 swap 写出针对 bt-fd 的运行；swap 写出是运行期间的全系统累计量。特征区是外扩后两块中较大的比例。}
\end{table}
\FloatBarrier
\section{FluxMG}
判据：双双精度原残差 $<10^{-8}$，时间含构建，单线程。多群算例是 Yue 等 2026 的类比，不是他们的矩阵。
\subsection{第 1 轮最终版本，测试集（对 BoomerAMG $\theta=0.25$，11 月 6 日）}
{\footnotesize\setlength{\tabcolsep}{4pt}
\begin{longtable}{llrrrrr}
\caption{测试集 19 类，SCORE 0.559，赢 17/19，零失败。AMG 取 3 遍最快，FluxMG 为单遍。来源 \texttt{fluxmg/final\_suite.txt}。}\\
\toprule
算例 & 类别 & Flux 迭代 & AMG 迭代 & Flux (s) & AMG (s) & 比值 \\
\midrule
\endfirsthead
\toprule
算例 & 类别 & Flux 迭代 & AMG 迭代 & Flux (s) & AMG (s) & 比值 \\
\midrule
\endhead
\bottomrule
\endlastfoot
\texttt{poisson:64} & 常系数 Poisson & 18 & 19 & 0.36 & 0.57 & \textbf{0.62} \\
\texttt{aniso:64:100} & 对齐各向异性，一个强方向 & 13 & 16 & 0.72 & 0.63 & 1.13 \\
\texttt{aniso:64:0.01} & 对齐各向异性，一个弱方向 & 18 & 17 & 0.44 & 0.61 & \textbf{0.73} \\
\texttt{lognorm3} & 高对比度，对数正态 & 18 & 19 & 0.40 & 0.60 & \textbf{0.66} \\
\texttt{tubes} & 高对比度，导电管道 & 24 & 20 & 0.49 & 0.70 & \textbf{0.69} \\
\texttt{insul} & 高对比度，绝缘夹杂 & 21 & 19 & 0.49 & 0.65 & \textbf{0.75} \\
\texttt{perc3} & 高对比度，渗流两相 & 23 & 19 & 0.65 & 0.73 & \textbf{0.90} \\
\texttt{g3t:64:2e-11:2e-4} & 3T，类 SC26\_01 & 22 & 30 & 1.24 & 2.55 & \textbf{0.49} \\
\texttt{g3t:64:2e-11:1} & 3T，强电子传导 & 34 & 30 & 2.03 & 3.22 & \textbf{0.63} \\
\texttt{g3t:64:1:1e-2} & 3T，强反馈 & 20 & 28 & 1.30 & 2.98 & \textbf{0.43} \\
\texttt{gen:64:1:8:0.5:0:2e-2:1} & 3T，扁方块 & 22 & 30 & 1.19 & 2.70 & \textbf{0.44} \\
\texttt{gen:64:1:2:1:4:2e-2:1} & 3T，渐变过渡 & 37 & 31 & 1.67 & 2.63 & \textbf{0.63} \\
\texttt{gen:64:3:2:1:0:2e-2:1} & 3T，3 胶囊 & 29 & 30 & 1.41 & 2.66 & \textbf{0.53} \\
\texttt{gen:64:10:2:1:0:2e-2:1} & 3T，10 胶囊 & 36 & 30 & 1.66 & 2.75 & \textbf{0.61} \\
\texttt{gen:64:10:2:1:0:2e-2:2} & 3T，10 胶囊，另一种子 & 36 & 29 & 1.68 & 2.68 & \textbf{0.63} \\
\texttt{gen:64:10:8:0.5:2:2e-2:3} & 3T，10 渐变扁方块 & 68 & 32 & 2.75 & 2.90 & \textbf{0.95} \\
\texttt{mg:32:8:1e-2:1e2:0:7} & 多群，AMG 易（类 U13-2） & 22 & 16 & 0.80 & 0.69 & 1.16 \\
\texttt{mg:32:4:1:1:0:7} & 多群，中（类 U92-1） & 25 & 25 & 0.33 & 0.51 & \textbf{0.66} \\
\texttt{mg:32:8:1e2:1e-2:1:7} & 多群，AMG 难（类 U107-4） & 121 & 109 & 1.82 & 2.47 & \textbf{0.74} \\
\end{longtable}}
\subsection{第 1 轮留出集：最终版本与起点版本}
{\footnotesize\setlength{\tabcolsep}{4pt}
\begin{longtable}{llrrrrrr}
\caption{最终版本（\texttt{best/}），留出 SCORE 0.665，赢 10/10。来源 \texttt{fluxmg/final\_holdout.txt}。}\\
\toprule
算例 & 类别 & Flux 迭代 & AMG 迭代 & Flux (s) & AMG (s) & 比值 & 区域数 $K$ \\
\midrule
\endfirsthead
\toprule
算例 & 类别 & Flux 迭代 & AMG 迭代 & Flux (s) & AMG (s) & 比值 & 区域数 $K$ \\
\midrule
\endhead
\bottomrule
\endlastfoot
\texttt{real} & 真实 SC26\_01 & 59 & 34 & 26.67 & 28.82 & \textbf{0.93} & 122 \\
\texttt{gen:64:5:2:1:0:2e-2:4} & 3T，5 胶囊，新种子 & 32 & 30 & 1.66 & 2.66 & \textbf{0.63} & 12 \\
\texttt{gen:64:10:2:1:0:2e-2:5} & 3T，10 胶囊，新种子 & 55 & 30 & 2.52 & 2.79 & \textbf{0.90} & 15 \\
\texttt{gen:64:1:2:2:0:2e-2:6} & 3T，长椭球 & 23 & 30 & 1.37 & 2.82 & \textbf{0.49} & 4 \\
\texttt{gen:64:2:1:1:1:1e-1:7} & 3T，八面体，10\% 扰动 & 37 & 30 & 1.69 & 2.57 & \textbf{0.65} & 5 \\
\texttt{g3t:64:1e-4:1e-1} & 3T，新耦合参数 & 23 & 30 & 1.60 & 2.70 & \textbf{0.59} & 4 \\
\texttt{mg:32:4:1e2:1:1:8} & 多群，新参数 & 14 & 26 & 0.21 & 0.53 & \textbf{0.39} & 3 \\
\texttt{mg:32:8:1:1e-2:0:9} & 多群，新参数 & 289 & 286 & 4.93 & 6.49 & \textbf{0.76} & 22 \\
\texttt{aniso:64:1000} & 对齐各向异性，更强 & 11 & 18 & 0.61 & 0.66 & \textbf{0.93} & 1 \\
\texttt{gen:128:3:2:1:0:2e-2:1} & 3T，3 胶囊，$m=128$ & 34 & 33 & 17.69 & 29.25 & \textbf{0.60} & 8 \\
\texttt{gen:128:10:2:1:0:2e-2:2} & 3T，10 胶囊，$m=128$ & 29 & 33 & 15.31 & 27.42 & \textbf{0.56} & 22 \\
\end{longtable}}
{\footnotesize\setlength{\tabcolsep}{4pt}
\begin{longtable}{llrrrrrr}
\caption{起点版本（循环开始时的代码），留出 SCORE 0.556，赢 9/10。}\\
\toprule
算例 & 类别 & Flux 迭代 & AMG 迭代 & Flux (s) & AMG (s) & 比值 & 区域数 $K$ \\
\midrule
\endfirsthead
\toprule
算例 & 类别 & Flux 迭代 & AMG 迭代 & Flux (s) & AMG (s) & 比值 & 区域数 $K$ \\
\midrule
\endhead
\bottomrule
\endlastfoot
\texttt{real} & 真实 SC26\_01 & 110 & 34 & 60.67 & 28.82 & 2.11 & 127 \\
\texttt{gen:64:5:2:1:0:2e-2:4} & 3T，5 胶囊，新种子 & 28 & 30 & 1.94 & 2.66 & \textbf{0.73} & 17 \\
\texttt{gen:64:10:2:1:0:2e-2:5} & 3T，10 胶囊，新种子 & 45 & 30 & 2.79 & 2.79 & 1.00 & 26 \\
\texttt{gen:64:1:2:2:0:2e-2:6} & 3T，长椭球 & 27 & 30 & 1.86 & 2.82 & \textbf{0.66} & 5 \\
\texttt{gen:64:2:1:1:1:1e-1:7} & 3T，八面体，10\% 扰动 & 53 & 30 & 3.09 & 2.57 & 1.20 & 4 \\
\texttt{g3t:64:1e-4:1e-1} & 3T，新耦合参数 & 27 & 30 & 1.87 & 2.70 & \textbf{0.69} & 5 \\
\texttt{mg:32:4:1e2:1:1:8} & 多群，新参数 & 14 & 26 & 0.27 & 0.53 & \textbf{0.51} & 9 \\
\texttt{mg:32:8:1:1e-2:0:9} & 多群，新参数 & 36 & 286 & 1.44 & 6.49 & \textbf{0.22} & 21 \\
\texttt{aniso:64:1000} & 对齐各向异性，更强 & 9 & 18 & 0.51 & 0.66 & \textbf{0.77} & 1 \\
\texttt{gen:128:3:2:1:0:2e-2:1} & 3T，3 胶囊，$m=128$ & 34 & 33 & 20.11 & 29.25 & \textbf{0.69} & 10 \\
\texttt{gen:128:10:2:1:0:2e-2:2} & 3T，10 胶囊，$m=128$ & 42 & 33 & 24.30 & 27.42 & \textbf{0.89} & 29 \\
\end{longtable}}
\subsection{第 30 轮：多群分组修复}
{\footnotesize\setlength{\tabcolsep}{4pt}
\begin{longtable}{rlrrr}
\caption{测试集交替配对运行（桌面负载约 3）。中位数：修复前 0.494，修复后 0.565。来源 \texttt{fluxmg/round2\_grouping/paired\_*.txt}。}\\
\toprule
顺序 & 代码 & SCORE & 胜 & 失败 \\
\midrule
\endfirsthead
\toprule
顺序 & 代码 & SCORE & 胜 & 失败 \\
\midrule
\endhead
\bottomrule
\endlastfoot
1 & 修复前 & 0.5697 & 17/19 & 0 \\
2 & 修复后 & 0.5926 & 17/19 & 0 \\
3 & 修复后 & 0.5484 & 16/19 & 0 \\
4 & 修复前 & 0.4931 & 17/19 & 0 \\
5 & 修复后 & 0.5837 & 17/19 & 0 \\
6 & 修复前 & 0.4930 & 16/19 & 0 \\
7 & 修复前 & 0.4058 & 16/19 & 0 \\
8 & 修复后 & 0.4793 & 16/19 & 0 \\
\end{longtable}}
{\scriptsize\setlength{\tabcolsep}{3pt}
\begin{longtable}{lrrrrrrrrl}
\caption{分组留出集 14 个（规则实现前固定）。“两两”为原 AMFR 两两分组，“规则”为加探针后的分组；比值对 BoomerAMG $\theta=0.25$。两两 SCORE 0.404（12/14），规则 SCORE 1.088（14/14）。分组一栏中 e 表示可由对角解消去的场，\textbar{} 分隔各组。来源 \texttt{round2\_grouping/grouping.txt}。}\\
\toprule
算例 & AMG & (s) & 两两 & (s) & 比值 & 规则 & (s) & 比值 & 规则分组 \\
\midrule
\endfirsthead
\toprule
算例 & AMG & (s) & 两两 & (s) & 比值 & 规则 & (s) & 比值 & 规则分组 \\
\midrule
\endhead
\bottomrule
\endlastfoot
\texttt{mg:32:6:10:0.03:0:21} & 300 & 4.87 & 261 & 2.67 & 0.55 & 24 & 0.60 & \textbf{0.12} & \texttt{e6e0\textbar 12345} \\
\texttt{mg:32:6:10:3:1:21} & 22 & 0.74 & 29 & 0.49 & 0.66 & 29 & 0.52 & \textbf{0.70} & \texttt{e0\textbar 5\textbar 1\textbar 2\textbar 6\textbar 3\textbar 4} \\
\texttt{mg:32:12:3:0.03:1:22} & 86 & 3.55 & 143 & 3.77 & 1.06 & 29 & 2.73 & \textbf{0.77} & \texttt{\textbar 0123456789101112} \\
\texttt{mg:32:12:30:0.3:0:22} & 100 & 3.75 & 81 & 2.52 & 0.67 & 19 & 1.68 & \textbf{0.45} & \texttt{e12e0e1e2e3\textbar 4567891011} \\
\texttt{mg:48:8:0.1:0.01:0:23} & 118 & 10.40 & 121 & 8.23 & 0.79 & 27 & 4.11 & \textbf{0.40} & \texttt{e8\textbar 01234567} \\
\texttt{mg:48:8:10:1e-3:1:23} & 211 & 17.14 & 391 & 18.43 & 1.08 & 39 & 4.60 & \textbf{0.27} & \texttt{e0\textbar 12345678} \\
\texttt{mg:48:4:3:0.3:0:24} & 72 & 3.54 & 65 & 1.85 & 0.52 & 23 & 1.23 & \textbf{0.35} & \texttt{e4\textbar 0123} \\
\texttt{mg:32:8:0.1:1:0:25} & 36 & 1.52 & 38 & 1.03 & 0.68 & 38 & 1.06 & \textbf{0.69} & \texttt{e8\textbar 0\textbar 7\textbar 1\textbar 2\textbar 3\textbar 4\textbar 5\textbar 6} \\
\texttt{mg:64:12:10:0.05:0:26} & 142 & 42.80 & 157 & 42.01 & 0.98 & 30 & 20.59 & \textbf{0.48} & \texttt{e12e0e1\textbar 234567891011} \\
\texttt{mg:64:6:1:10:1:26} & 15 & 6.20 & 24 & 3.98 & 0.64 & 24 & 4.05 & \textbf{0.65} & \texttt{\textbar 0\textbar 5\textbar 1\textbar 6\textbar 2\textbar 3\textbar 4} \\
\texttt{blob:32:6:3:1e-2:0:31:0.1} & 89 & 1.56 & 87 & 1.21 & 0.77 & 22 & 0.67 & \textbf{0.43} & \texttt{e6\textbar 012345} \\
\texttt{blob:32:6:3:1e-2:1:31:0.1} & 14 & 0.55 & 21 & 0.49 & 0.89 & 21 & 0.47 & \textbf{0.85} & \texttt{\textbar 0\textbar 5\textbar 1\textbar 2\textbar 6\textbar 3\textbar 4} \\
\texttt{blob:32:12:30:1e-2:0:32:0.03} & 153 & 4.95 & 134 & 3.58 & 0.72 & 23 & 1.74 & \textbf{0.35} & \texttt{e12e0e1e2e3\textbar 4567891011} \\
\texttt{blob:48:8:1:1e-2:0:33:0.03} & 30 & 3.15 & 31 & 2.53 & 0.80 & 31 & 2.66 & \textbf{0.85} & \texttt{e8\textbar 0\textbar 1\textbar 7\textbar 2\textbar 3\textbar 4\textbar 5\textbar 6} \\
\end{longtable}}
{\footnotesize\setlength{\tabcolsep}{4pt}
\begin{longtable}{llrrrrr}
\caption{修复后的第 1 轮留出集，SCORE 0.812，赢 10/10（同一负载下）。来源 \texttt{round2\_grouping/holdout.txt}。}\\
\toprule
算例 & 类别 & Flux 迭代 & AMG 迭代 & Flux (s) & AMG (s) & 比值 \\
\midrule
\endfirsthead
\toprule
算例 & 类别 & Flux 迭代 & AMG 迭代 & Flux (s) & AMG (s) & 比值 \\
\midrule
\endhead
\bottomrule
\endlastfoot
\texttt{gen:64:5:2:1:0:2e-2:4} & 3T，5 胶囊，新种子 & 32 & 30 & 1.73 & 2.66 & \textbf{0.65} \\
\texttt{gen:64:10:2:1:0:2e-2:5} & 3T，10 胶囊，新种子 & 55 & 30 & 2.54 & 2.79 & \textbf{0.91} \\
\texttt{gen:64:1:2:2:0:2e-2:6} & 3T，长椭球 & 23 & 30 & 1.33 & 2.82 & \textbf{0.47} \\
\texttt{gen:64:2:1:1:1:1e-1:7} & 3T，八面体，10\% 扰动 & 37 & 30 & 1.82 & 2.57 & \textbf{0.71} \\
\texttt{g3t:64:1e-4:1e-1} & 3T，新耦合参数 & 23 & 30 & 1.76 & 2.70 & \textbf{0.65} \\
\texttt{mg:32:4:1e2:1:1:8} & 多群，新参数 & 14 & 26 & 0.22 & 0.53 & \textbf{0.42} \\
\texttt{mg:32:8:1:1e-2:0:9} & 多群，新参数 & 26 & 286 & 1.30 & 6.49 & \textbf{0.20} \\
\texttt{aniso:64:1000} & 对齐各向异性，更强 & 11 & 18 & 0.59 & 0.66 & \textbf{0.90} \\
\texttt{gen:128:3:2:1:0:2e-2:1} & 3T，3 胶囊，$m=128$ & 34 & 33 & 20.40 & 29.25 & \textbf{0.70} \\
\texttt{gen:128:10:2:1:0:2e-2:2} & 3T，10 胶囊，$m=128$ & 29 & 33 & 14.48 & 27.42 & \textbf{0.53} \\
\end{longtable}}
{\footnotesize\setlength{\tabcolsep}{4pt}
\begin{longtable}{rlrrlrrr}
\caption{修复后真实 SC26\_01，3 次独立运行。来源 \texttt{round2\_grouping/real.txt}。}\\
\toprule
运行 & 分组 & 探针 & 总迭代 & 各次修正迭代 & 最终相对残差 & 构建 (s) & 总时间 (s) \\
\midrule
\endfirsthead
\toprule
运行 & 分组 & 探针 & 总迭代 & 各次修正迭代 & 最终相对残差 & 构建 (s) & 总时间 (s) \\
\midrule
\endhead
\bottomrule
\endlastfoot
1 & \texttt{e1\textbar 2\textbar 0} & 3.13e-05 & 59 & 26, 14, 19 & 7.90e-15 & 3.68 & 23.59 \\
2 & \texttt{e1\textbar 2\textbar 0} & 3.13e-05 & 59 & 26, 14, 19 & 7.90e-15 & 3.57 & 23.63 \\
3 & \texttt{e1\textbar 2\textbar 0} & 3.13e-05 & 59 & 26, 14, 19 & 7.90e-15 & 3.57 & 24.51 \\
\end{longtable}}
\begin{table}[H]
\centering\small
\begin{tabular}{lrrrrr}
\toprule
独立进程，各 3 遍 & 方法 & 迭代 & 时间 (s) & 配对比值 & 进程峰值内存 (GB) \\
\midrule
真实 SC26\_01 & FluxMG & 59 & 22.9，24.0，26.0 & 0.81，0.81，0.87 & 5.0--5.9 \\
 & BoomerAMG & 34 & 28.3，29.8，30.0 & & 5.2--5.4 \\
10 胶囊，$m=128$ & FluxMG & 41 & 19.8，19.7，19.2 & 0.74，0.73，0.73 & 5.0--5.7 \\
 & BoomerAMG & 33 & 26.7，27.0，26.4 & & 6.0--6.5 \\
\bottomrule
\end{tabular}
\caption{FluxMG 迭代更多，但每步便宜，构建约 3.5\,s（部分仍是 Python）。内存是整个进程的峰值，含公共的矩阵和 DD 表，以及 FGMRES 比 GMRES 多存的向量。}
\end{table}
\subsection{11 月 9 日预注册比较：FluxMG、hypre 结构多重网格和 BoomerAMG}
每格为“时间 (s)/迭代”，2 次打乱顺序重复的中位数；任一次失败即记为失败。$^\dagger$ 表示结构方法内部有一组多场用一次 BoomerAMG V-cycle（混合）。AMG 列为 $\theta=0.25$；“AMG 最优”取 D.5 中每个算例最快的设置。比值 = FluxMG / 对照。来源 \texttt{fluxmg/structured/results\_20261109.jsonl}、\texttt{fluxmg/amg\_configs.jsonl}。
{\scriptsize\setlength{\tabcolsep}{3pt}
\begin{longtable}{lrrrrrlrr}
\caption{测试集。}\\
\toprule
算例 & FluxMG & PFMG & PFMG-rb & SMG & AMG .25 & AMG 最优（设置） & /结构 & /AMG \\
\midrule
\endfirsthead
\toprule
算例 & FluxMG & PFMG & PFMG-rb & SMG & AMG .25 & AMG 最优（设置） & /结构 & /AMG \\
\midrule
\endhead
\bottomrule
\endlastfoot
\texttt{poisson:64} & 0.38/18 & 0.23/15 & 0.19/19 & 0.65/8 & 0.59/19 & 0.58/19 fast.25 & 1.96 & 0.65 \\
\texttt{aniso:64:100} & 0.73/13 & 0.29/13 & 0.21/16 & 0.64/8 & 0.64/16 & 0.60/16 fast.5 & 3.41 & 1.22 \\
\texttt{aniso:64:0.01} & 0.45/18 & 0.28/15 & 0.26/26 & 0.63/8 & 0.63/17 & 0.59/10 falgout.25 & 1.76 & 0.76 \\
\texttt{lognorm3} & 0.39/18 & 0.26/17 & 0.20/18 & 0.77/10 & 0.62/19 & 0.60/19 fast.25 & 1.90 & 0.65 \\
\texttt{tubes} & 0.51/24 & 0.84/76 & 0.23/22 & 0.79/11 & 0.73/20 & 0.69/20 fast.25 & 2.24 & 0.74 \\
\texttt{insul} & 0.47/21 & 0.37/32 & 0.24/22 & 0.71/9 & 0.65/19 & 0.62/19 fast.25 & 1.98 & 0.76 \\
\texttt{perc3} & 0.72/23 & 0.49/44 & 0.24/26 & 1.00/14 & 0.72/19 & 0.71/20 fast.5 & 3.00 & 1.02 \\
\texttt{g3t:64:2e-11:2e-4} & 1.26/22 & 1.66/45 & 失败 & 19.09/141 & 2.65/30 & 2.55/30 fast.25 & 0.76 & 0.50 \\
\texttt{g3t:64:2e-11:1} & 2.07/33 & 2.23/23$^\dagger$ & 2.59/31$^\dagger$ & 3.15/21$^\dagger$ & 3.31/30 & 3.21/30 fast.25 & 0.93 & 0.65 \\
\texttt{g3t:64:1:1e-2} & 1.44/20 & 1.95/20$^\dagger$ & 1.93/20$^\dagger$ & 1.92/20$^\dagger$ & 2.94/28 & 2.35/30 sys.u & 0.75 & 0.61 \\
\texttt{gen:64:1:8:0.5:0:2e-2:1} & 1.28/22 & 1.20/29 & 失败 & 9.32/74 & 2.76/30 & 2.61/30 fast.25 & 1.06 & 0.49 \\
\texttt{gen:64:1:2:1:4:2e-2:1} & 1.82/37 & 1.41/37 & 1.98/69 & 10.29/79 & 2.62/31 & 2.65/31 fast.25 & 1.29 & 0.69 \\
\texttt{gen:64:3:2:1:0:2e-2:1} & 1.52/30 & 失败 & 失败 & 失败 & 2.63/30 & 2.56/30 fast.25 & 结构均失败 & 0.59 \\
\texttt{gen:64:10:2:1:0:2e-2:1} & 1.78/36 & 失败 & 失败 & 失败 & 2.75/30 & 2.71/30 fast.25 & 结构均失败 & 0.66 \\
\texttt{gen:64:10:2:1:0:2e-2:2} & 1.80/36 & 失败 & 失败 & 失败 & 2.74/29 & 2.79/29 fast.25 & 结构均失败 & 0.64 \\
\texttt{gen:64:10:8:0.5:2:2e-2:3} & 3.07/67 & 失败 & 失败 & 失败 & 2.87/32 & 2.87/32 fast.25 & 结构均失败 & 1.07 \\
\texttt{mg:32:8:1e-2:1e2:0:7} & 0.82/22 & 0.58/15 & 0.55/18 & 1.16/9 & 0.75/16 & 0.52/16 sys.n & 1.48 & 1.59 \\
\texttt{mg:32:4:1:1:0:7} & 0.33/25 & 0.27/23 & 0.25/24 & 1.00/23 & 0.52/25 & 0.32/25 sys.n & 1.32 & 1.04 \\
\texttt{mg:32:8:1e2:1e-2:1:7} & 0.83/23 & 1.51/41$^\dagger$ & 1.49/41$^\dagger$ & 1.53/41$^\dagger$ & 2.58/109 & 1.56/109 sys.n & 0.56 & 0.53 \\
\end{longtable}}
{\scriptsize\setlength{\tabcolsep}{3pt}
\begin{longtable}{lrrrrrlrr}
\caption{留出集。}\\
\toprule
算例 & FluxMG & PFMG & PFMG-rb & SMG & AMG .25 & AMG 最优（设置） & /结构 & /AMG \\
\midrule
\endfirsthead
\toprule
算例 & FluxMG & PFMG & PFMG-rb & SMG & AMG .25 & AMG 最优（设置） & /结构 & /AMG \\
\midrule
\endhead
\bottomrule
\endlastfoot
\texttt{gen:64:5:2:1:0:2e-2:4} & 1.85/32 & 失败 & 失败 & 失败 & 2.67/30 & 2.63/30 fast.25 & 结构均失败 & 0.70 \\
\texttt{gen:64:10:2:1:0:2e-2:5} & 2.33/55 & 失败 & 失败 & 失败 & 2.94/30 & 2.72/30 fast.25 & 结构均失败 & 0.86 \\
\texttt{gen:64:1:2:2:0:2e-2:6} & 1.33/23 & 1.80/50 & 失败 & 29.22/229 & 2.83/30 & 2.83/30 fast.25 & 0.74 & 0.47 \\
\texttt{gen:64:2:1:1:1:1e-1:7} & 1.74/37 & 14.42/600 & 失败 & 43.96/346 & 2.59/30 & 2.50/30 fast.25 & 0.12 & 0.70 \\
\texttt{g3t:64:1e-4:1e-1} & 1.68/23 & 1.90/23$^\dagger$ & 2.36/31$^\dagger$ & 3.06/21$^\dagger$ & 2.71/30 & 2.54/30 fast.25 & 0.89 & 0.66 \\
\texttt{mg:32:4:1e2:1:1:8} & 0.21/14 & 0.25/12$^\dagger$ & 0.24/11$^\dagger$ & 0.34/11$^\dagger$ & 0.54/26 & 0.31/26 sys.n & 0.90 & 0.68 \\
\texttt{mg:32:8:1:1e-2:0:9} & 1.34/26 & 6.47/152$^\dagger$ & 6.55/152$^\dagger$ & 6.47/152$^\dagger$ & 6.43/286 & 4.69/285 sys.n & 0.21 & 0.29 \\
\texttt{aniso:64:1000} & 0.59/11 & 0.34/16 & 0.24/20 & 0.74/10 & 0.64/18 & 0.52/8 falgout.25 & 2.47 & 1.15 \\
\texttt{gen:128:3:2:1:0:2e-2:1} & 17.49/34 & 181.74/600 & 失败 & 失败 & 25.79/33 & 25.56/33 fast.25 & 0.10 & 0.68 \\
\texttt{gen:128:10:2:1:0:2e-2:2} & 15.15/29 & 失败 & 失败 & 失败 & 28.82/33 & 26.92/33 fast.25 & 结构均失败 & 0.56 \\
\end{longtable}}
{\scriptsize\setlength{\tabcolsep}{3pt}
\begin{longtable}{lrrrrrlrr}
\caption{分组留出集。}\\
\toprule
算例 & FluxMG & PFMG & PFMG-rb & SMG & AMG .25 & AMG 最优（设置） & /结构 & /AMG \\
\midrule
\endfirsthead
\toprule
算例 & FluxMG & PFMG & PFMG-rb & SMG & AMG .25 & AMG 最优（设置） & /结构 & /AMG \\
\midrule
\endhead
\bottomrule
\endlastfoot
\texttt{mg:32:6:10:0.03:0:21} & 0.61/24 & 3.65/144$^\dagger$ & 3.55/144$^\dagger$ & 3.53/144$^\dagger$ & 4.79/300 & 3.30/300 sys.n & 0.17 & 0.18 \\
\texttt{mg:32:6:10:3:1:21} & 0.52/29 & 0.49/30 & 0.45/30 & 1.82/29 & 0.71/22 & 0.52/20 sys.u & 1.16 & 0.99 \\
\texttt{mg:32:12:3:0.03:1:22} & 2.69/29 & 4.12/86$^\dagger$ & 4.12/86$^\dagger$ & 4.20/86$^\dagger$ & 3.46/86 & 2.70/74 sys.u & 0.65 & 0.99 \\
\texttt{mg:32:12:30:0.3:0:22} & 1.55/19 & 1.61/13$^\dagger$ & 1.72/13$^\dagger$ & 1.70/13$^\dagger$ & 3.75/100 & 2.17/100 sys.n & 0.96 & 0.71 \\
\texttt{mg:48:8:0.1:0.01:0:23} & 4.08/27 & 13.36/63$^\dagger$ & 13.03/63$^\dagger$ & 12.84/63$^\dagger$ & 10.15/118 & 7.59/118 sys.n & 0.32 & 0.54 \\
\texttt{mg:48:8:10:1e-3:1:23} & 4.66/39 & 19.04/220$^\dagger$ & 19.21/220$^\dagger$ & 18.87/220$^\dagger$ & 17.66/211 & 13.73/182 sys.u & 0.25 & 0.34 \\
\texttt{mg:48:4:3:0.3:0:24} & 1.31/23 & 2.57/38$^\dagger$ & 2.67/38$^\dagger$ & 2.65/38$^\dagger$ & 3.52/72 & 2.53/72 sys.u & 0.51 & 0.52 \\
\texttt{mg:32:8:0.1:1:0:25} & 1.09/38 & 0.89/38 & 0.81/39 & 3.22/38 & 1.55/36 & 0.79/36 sys.n & 1.35 & 1.38 \\
\texttt{mg:64:12:10:0.05:0:26} & 21.54/30 & 58.62/83$^\dagger$ & 60.07/83$^\dagger$ & 66.35/83$^\dagger$ & 42.75/142 & 29.02/142 sys.n & 0.37 & 0.74 \\
\texttt{mg:64:6:1:10:1:26} & 4.20/24 & 3.39/17 & 2.90/18 & 8.82/16 & 6.89/15 & 5.17/15 sys.u & 1.45 & 0.81 \\
\texttt{blob:32:6:3:1e-2:0:31:0.1} & 0.71/22 & 1.54/49$^\dagger$ & 1.48/49$^\dagger$ & 1.49/49$^\dagger$ & 1.62/89 & 1.14/89 sys.n & 0.48 & 0.62 \\
\texttt{blob:32:6:3:1e-2:1:31:0.1} & 0.52/21 & 0.38/16 & 0.33/15 & 1.04/12 & 0.57/14 & 0.51/15 sys.u & 1.56 & 1.02 \\
\texttt{blob:32:12:30:1e-2:0:32:0.03} & 1.68/23 & 1.91/22$^\dagger$ & 1.85/22$^\dagger$ & 1.93/22$^\dagger$ & 5.14/153 & 2.91/153 sys.n & 0.91 & 0.58 \\
\texttt{blob:48:8:1:1e-2:0:33:0.03} & 2.69/31 & 2.53/31 & 2.21/31 & 8.12/30 & 3.46/30 & 2.51/29 sys.n & 1.22 & 1.07 \\
\end{longtable}}
{\scriptsize\setlength{\tabcolsep}{3pt}
\begin{longtable}{lrrrrrlrr}
\caption{真实矩阵。}\\
\toprule
算例 & FluxMG & PFMG & PFMG-rb & SMG & AMG .25 & AMG 最优（设置） & /结构 & /AMG \\
\midrule
\endfirsthead
\toprule
算例 & FluxMG & PFMG & PFMG-rb & SMG & AMG .25 & AMG 最优（设置） & /结构 & /AMG \\
\midrule
\endhead
\bottomrule
\endlastfoot
\texttt{real} & 23.16/59 & 失败 & 失败 & 失败 & 27.94/34 & 29.01/34 fast.25 & 结构均失败 & 0.80 \\
\end{longtable}}
预注册预测 P5（FluxMG 在 $\ge70\%$ 的测试和留出算例上快于最好的结构方法）：59\%，不成立。
\subsection{BoomerAMG 全部设置}
每格为“时间 (s)/迭代”，同一 \texttt{solve\_ir} 路径；重复运行取中位数。粗体为该算例最快的设置。fast = HMIS 粗化 + 激进粗化，falgout = Falgout 粗化，noagg = 不用激进粗化，数字为强连接阈值 $\theta$；sys.u/sys.n 为多场 systems 设置（未知量方式 / 节点方式），只对多场算例运行。来源 \texttt{fluxmg/amg\_configs.jsonl}。
{\tiny\setlength{\tabcolsep}{2.5pt}
\begin{longtable}{lrrrrrrrrrr}
\caption{测试集。}\\
\toprule
算例 & fast.25 & fast.5 & fast.7 & fast.9 & falgout.25 & falgout.5 & noagg.25 & noagg.5 & sys.u & sys.n \\
\midrule
\endfirsthead
\toprule
算例 & fast.25 & fast.5 & fast.7 & fast.9 & falgout.25 & falgout.5 & noagg.25 & noagg.5 & sys.u & sys.n \\
\midrule
\endhead
\bottomrule
\endlastfoot
\texttt{poisson:64} & \textbf{0.58/19} & 0.60/19 & 0.63/21 & 0.63/22 & 0.93/8 & 0.69/9 & 1.11/15 & 1.18/15 & -- & -- \\
\texttt{aniso:64:100} & 0.62/16 & \textbf{0.60/16} & 0.60/16 & 0.60/16 & 0.62/8 & 0.61/8 & 0.75/13 & 0.62/10 & -- & -- \\
\texttt{aniso:64:0.01} & 0.60/17 & 0.64/17 & 0.65/18 & 0.66/19 & \textbf{0.59/10} & 1.17/12 & 0.76/11 & 1.09/11 & -- & -- \\
\texttt{lognorm3} & \textbf{0.60/19} & 0.69/22 & 0.85/24 & 1.30/35 & 1.12/10 & 1.73/13 & 1.15/14 & 1.20/16 & -- & -- \\
\texttt{tubes} & \textbf{0.69/20} & 0.71/22 & 0.75/24 & 0.78/26 & 1.51/12 & 1.37/12 & 1.06/12 & 0.99/14 & -- & -- \\
\texttt{insul} & \textbf{0.62/19} & 0.69/21 & 0.67/21 & 0.74/25 & 0.90/10 & 1.31/12 & 1.17/15 & 1.03/15 & -- & -- \\
\texttt{perc3} & 0.71/19 & \textbf{0.71/20} & 0.76/24 & 0.81/28 & 1.37/10 & 1.22/12 & 1.05/13 & 0.94/15 & -- & -- \\
\texttt{g3t:64:2e-11:2e-4} & \textbf{2.55/30} & 2.66/31 & 2.87/35 & 3.23/42 & 2.94/14 & 3.01/17 & 3.33/23 & 3.84/23 & 2.80/40 & 6.00/79 \\
\texttt{g3t:64:2e-11:1} & \textbf{3.21/30} & 3.47/32 & 3.61/36 & 4.04/44 & 5.44/17 & 5.38/18 & 5.03/22 & 6.43/24 & 7.06/97 & 5.75/77 \\
\texttt{g3t:64:1:1e-2} & 2.86/28 & 3.11/29 & 3.21/34 & 3.61/42 & 4.93/14 & 4.12/17 & 3.98/21 & 5.35/21 & \textbf{2.35/30} & 2.95/30 \\
\texttt{gen:64:1:8:0.5:0:2e-2:1} & \textbf{2.61/30} & 2.93/33 & 2.83/35 & 3.26/43 & 3.66/14 & 4.19/19 & 2.97/17 & 3.98/27 & 2.94/42 & 4.99/62 \\
\texttt{gen:64:1:2:1:4:2e-2:1} & \textbf{2.65/31} & 2.78/33 & 2.99/37 & 3.55/46 & 3.17/15 & 4.62/21 & 3.69/24 & 4.00/27 & 5.53/97 & 7.12/94 \\
\texttt{gen:64:3:2:1:0:2e-2:1} & \textbf{2.56/30} & 2.81/31 & 3.12/38 & 3.68/49 & 3.12/14 & 4.18/19 & 3.61/23 & 3.88/26 & 2.86/43 & 4.76/60 \\
\texttt{gen:64:10:2:1:0:2e-2:1} & \textbf{2.71/30} & 2.87/32 & 3.36/41 & 3.87/50 & 3.37/15 & 4.20/20 & 3.62/23 & 3.85/26 & 2.86/42 & 4.86/62 \\
\texttt{gen:64:10:2:1:0:2e-2:2} & \textbf{2.79/29} & 2.84/32 & 3.40/42 & 4.28/56 & 3.70/20 & 4.85/27 & 3.80/26 & 4.16/29 & 3.01/43 & 4.69/59 \\
\texttt{gen:64:10:8:0.5:2:2e-2:3} & \textbf{2.87/32} & 2.91/34 & 3.35/40 & 4.19/54 & 3.60/16 & 4.61/22 & 3.68/25 & 4.01/28 & 5.62/95 & 7.33/97 \\
\texttt{mg:32:8:1e-2:1e2:0:7} & 0.73/16 & 0.71/16 & 0.75/18 & 0.93/25 & 1.22/8 & 1.65/11 & 1.18/11 & 1.19/13 & 0.65/16 & \textbf{0.52/16} \\
\texttt{mg:32:4:1:1:0:7} & 0.50/25 & 0.54/25 & 0.56/25 & 0.56/26 & 1.86/25 & 1.96/25 & 0.88/25 & 0.92/25 & 0.34/25 & \textbf{0.32/25} \\
\texttt{mg:32:8:1e2:1e-2:1:7} & 2.43/109 & 2.52/109 & 2.39/109 & 2.44/109 & 2.70/109 & 2.77/109 & 2.60/109 & 2.68/109 & 1.76/91 & \textbf{1.56/109} \\
\end{longtable}}
{\tiny\setlength{\tabcolsep}{2.5pt}
\begin{longtable}{lrrrrrrrrrr}
\caption{留出集。}\\
\toprule
算例 & fast.25 & fast.5 & fast.7 & fast.9 & falgout.25 & falgout.5 & noagg.25 & noagg.5 & sys.u & sys.n \\
\midrule
\endfirsthead
\toprule
算例 & fast.25 & fast.5 & fast.7 & fast.9 & falgout.25 & falgout.5 & noagg.25 & noagg.5 & sys.u & sys.n \\
\midrule
\endhead
\bottomrule
\endlastfoot
\texttt{gen:64:5:2:1:0:2e-2:4} & \textbf{2.63/30} & 2.80/32 & 3.13/38 & 3.57/46 & 3.05/14 & 4.14/19 & 3.59/23 & 3.93/26 & 2.79/42 & 4.83/62 \\
\texttt{gen:64:10:2:1:0:2e-2:5} & \textbf{2.72/30} & 2.96/33 & 3.36/42 & 3.94/51 & 3.25/15 & 4.25/20 & 3.63/24 & 3.94/27 & 2.96/43 & 4.95/62 \\
\texttt{gen:64:1:2:2:0:2e-2:6} & \textbf{2.83/30} & 2.90/32 & 3.02/36 & 3.54/46 & 3.39/14 & 4.34/20 & 3.70/24 & 3.88/26 & 3.07/43 & 5.85/75 \\
\texttt{gen:64:2:1:1:1:1e-1:7} & \textbf{2.50/30} & 2.78/32 & 3.03/36 & 3.55/46 & 3.53/15 & 4.35/19 & 3.60/23 & 4.01/27 & 3.94/64 & 4.83/64 \\
\texttt{g3t:64:1e-4:1e-1} & \textbf{2.54/30} & 2.77/32 & 2.98/36 & 3.40/43 & 2.74/11 & 3.09/18 & 3.52/23 & 3.76/22 & 5.83/99 & 5.99/80 \\
\texttt{mg:32:4:1e2:1:1:8} & 0.52/26 & 0.57/26 & 0.55/26 & 0.55/26 & 1.27/26 & 0.78/26 & 0.69/26 & 0.63/26 & 0.32/24 & \textbf{0.31/26} \\
\texttt{mg:32:8:1:1e-2:0:9} & 6.33/286 & 6.51/286 & 6.59/286 & 6.86/285 & 9.01/285 & 12.05/285 & 8.15/285 & 8.60/285 & 5.10/285 & \textbf{4.69/285} \\
\texttt{aniso:64:1000} & 0.63/18 & 0.62/18 & 0.63/18 & 0.63/18 & \textbf{0.52/8} & 0.56/9 & 0.52/8 & 0.54/9 & -- & -- \\
\texttt{gen:128:3:2:1:0:2e-2:1} & \textbf{25.56/33} & 28.27/36 & 31.23/43 & 38.07/57 & 39.05/15 & 47.73/22 & 38.18/27 & 46.43/35 & 29.37/45 & 77.92/106 \\
\texttt{gen:128:10:2:1:0:2e-2:2} & \textbf{26.92/33} & 31.17/37 & 35.90/48 & 41.99/63 & 39.45/15 & 44.73/22 & 38.59/26 & 44.75/35 & 29.00/45 & 48.64/64 \\
\end{longtable}}
{\tiny\setlength{\tabcolsep}{2.5pt}
\begin{longtable}{lrrrrrrrrrr}
\caption{分组留出集。}\\
\toprule
算例 & fast.25 & fast.5 & fast.7 & fast.9 & falgout.25 & falgout.5 & noagg.25 & noagg.5 & sys.u & sys.n \\
\midrule
\endfirsthead
\toprule
算例 & fast.25 & fast.5 & fast.7 & fast.9 & falgout.25 & falgout.5 & noagg.25 & noagg.5 & sys.u & sys.n \\
\midrule
\endhead
\bottomrule
\endlastfoot
\texttt{mg:32:6:10:0.03:0:21} & 4.78/300 & 4.69/300 & 4.66/300 & 4.73/300 & 5.83/300 & 6.85/300 & 5.48/300 & 5.57/300 & 3.79/300 & \textbf{3.30/300} \\
\texttt{mg:32:6:10:3:1:21} & 0.70/22 & 0.74/22 & 0.73/22 & 0.73/22 & 2.72/22 & 1.40/22 & 1.07/22 & 0.99/22 & \textbf{0.52/20} & 0.55/22 \\
\texttt{mg:32:12:3:0.03:1:22} & 3.37/86 & 3.35/86 & 3.36/86 & 3.49/86 & 4.31/85 & 5.28/85 & 4.16/86 & 4.26/86 & \textbf{2.70/74} & 3.06/86 \\
\texttt{mg:32:12:30:0.3:0:22} & 3.60/100 & 3.52/100 & 3.51/100 & 3.43/100 & 5.05/100 & 4.58/100 & 4.18/100 & 3.91/100 & 2.60/100 & \textbf{2.17/100} \\
\texttt{mg:48:8:0.1:0.01:0:23} & 9.90/118 & 10.52/118 & 10.51/118 & 11.64/118 & 18.82/117 & 26.02/118 & 14.97/118 & 16.15/118 & 9.16/118 & \textbf{7.59/118} \\
\texttt{mg:48:8:10:1e-3:1:23} & 16.65/211 & 16.84/211 & 16.68/208 & 17.57/213 & 21.68/212 & 26.16/209 & 20.20/211 & 20.96/212 & \textbf{13.73/182} & 15.53/209 \\
\texttt{mg:48:4:3:0.3:0:24} & 3.46/72 & 3.84/72 & 3.80/72 & 3.79/72 & 9.83/71 & 10.38/71 & 5.04/72 & 5.51/72 & \textbf{2.53/72} & 2.57/72 \\
\texttt{mg:32:8:0.1:1:0:25} & 1.52/36 & 1.64/36 & 1.34/36 & 1.49/37 & 11.37/36 & 10.86/36 & 3.18/36 & 3.46/36 & 0.96/36 & \textbf{0.79/36} \\
\texttt{mg:64:12:10:0.05:0:26} & 42.09/142 & 42.25/142 & 42.11/142 & 43.59/142 & 56.57/142 & 63.63/142 & 50.17/142 & 52.47/142 & 34.34/142 & \textbf{29.02/142} \\
\texttt{mg:64:6:1:10:1:26} & 6.09/15 & 6.01/15 & 6.30/17 & 7.62/23 & 11.93/10 & 16.12/11 & 9.27/12 & 8.43/12 & \textbf{5.17/15} & 5.24/17 \\
\texttt{blob:32:6:3:1e-2:0:31:0.1} & 1.57/89 & 1.58/89 & 1.62/89 & 1.66/89 & 2.16/89 & 2.80/89 & 2.01/89 & 2.06/89 & 1.26/89 & \textbf{1.14/89} \\
\texttt{blob:32:6:3:1e-2:1:31:0.1} & 0.57/14 & 0.60/16 & 0.58/16 & 0.69/20 & 1.25/7 & 1.21/8 & 0.86/11 & 0.85/12 & \textbf{0.51/15} & 0.57/19 \\
\texttt{blob:32:12:30:1e-2:0:32:0.03} & 4.95/153 & 4.95/153 & 4.89/153 & 4.87/153 & 6.06/153 & 6.25/153 & 5.56/153 & 5.49/153 & 3.70/153 & \textbf{2.91/153} \\
\texttt{blob:48:8:1:1e-2:0:33:0.03} & 3.10/30 & 3.25/30 & 3.13/30 & 3.39/30 & 4.75/29 & 6.32/30 & 4.70/30 & 4.71/30 & 2.71/30 & \textbf{2.51/29} \\
\end{longtable}}
{\tiny\setlength{\tabcolsep}{2.5pt}
\begin{longtable}{lrrrrrrrrrr}
\caption{真实矩阵。}\\
\toprule
算例 & fast.25 & fast.5 & fast.7 & fast.9 & falgout.25 & falgout.5 & noagg.25 & noagg.5 & sys.u & sys.n \\
\midrule
\endfirsthead
\toprule
算例 & fast.25 & fast.5 & fast.7 & fast.9 & falgout.25 & falgout.5 & noagg.25 & noagg.5 & sys.u & sys.n \\
\midrule
\endhead
\bottomrule
\endlastfoot
\texttt{real} & \textbf{29.01/34} & 29.80/37 & 32.43/44 & 44.80/67 & 47.23/16 & 42.71/23 & 33.47/26 & 47.34/31 & 225.71/450 & 失败 \\
\end{longtable}}
\subsection{autoresearch 迭代记录}
{\scriptsize\setlength{\tabcolsep}{3pt}
\begin{longtable}{rlll}
\caption{迭代记录原文（英文）。第 14--29 轮为 16 组同批参数扫描，合并为一行。第 31--33 轮在 11 月 9 日，未按 SCORE 判定。来源 \texttt{fluxmg/autoresearch.jsonl}。}\\
\toprule
轮 & SCORE & 结果 & 改动 \\
\midrule
\endfirsthead
\toprule
轮 & SCORE & 结果 & 改动 \\
\midrule
\endhead
\bottomrule
\endlastfoot
0 & -0.0082 & baseline & \parbox[t]{9.2cm}{\raggedright Prototype: flux form, K-cycle, lines, semi-coarsening, region vectors} \\
1 & -0.0899 & discard & \parbox[t]{9.2cm}{\raggedright Line trigger compares with the weakest other direction (in-plane strength)} \\
2 & 0.1526 & keep & \parbox[t]{9.2cm}{\raggedright Region vectors on every coarse level, correction after each coarse solve; graded boxes converge (70 its)} \\
3 & 0.2801 & keep & \parbox[t]{9.2cm}{\raggedright Automatic field reduction (AMFR rules) with one Core cycle per group; multigroup hard worsens 23 -\textgreater{} 121 its} \\
4 & 0.2352 & discard & \parbox[t]{9.2cm}{\raggedright TAU = 10. Real matrix reaches parity with BoomerAMG but the suite score does not improve; strong candidate to retest} \\
5 & 0.2129 & discard & \parbox[t]{9.2cm}{\raggedright AMFR spectral grouping} \\
6 & 0.0895 & discard & \parbox[t]{9.2cm}{\raggedright Convergence probe (8 cycles, limit 0.6) chooses reduced or monolithic correctly, but its cost lowers the score} \\
7 & 0.3225 & keep & \parbox[t]{9.2cm}{\raggedright Lines only where 4 consecutive anisotropic cells exist: isolated jumps no longer trigger them; insul and perc3 flip to wins} \\
8 & 0.3276 & discard & \parbox[t]{9.2cm}{\raggedright TAU = 10 on the new base; no score gain, and LINE\_RUN = 4 is seen to hurt the real matrix} \\
9 & 0.2241 & discard & \parbox[t]{9.2cm}{\raggedright TAU = 10, LINE\_RUN = 1, ANISO = 30: real matrix wins clearly, but whole-line solves make insul and perc3 lose} \\
10 & 0.3166 & hold & \parbox[t]{9.2cm}{\raggedright Iteration 9 settings with line solves restricted to anisotropic segments: suite score unchanged, real matrix now wins. Held in candidate\_iter10/ pending the user's decision on tie-breaks; score includes the real case} \\
10 & 0.3166 & keep & \parbox[t]{9.2cm}{\raggedright Held result of iteration 10 kept under the new tie-break} \\
11 & 0.3301 & discard & \parbox[t]{9.2cm}{\raggedright Per-axis region cut plus no lines on semi-coarsened axes: the real matrix collapses; no suite gain} \\
12 & 0.3022 & keep & \parbox[t]{9.2cm}{\raggedright Bug fix outside the score rule: coarsening stopped once an axis reached one layer, leaving 4096 cells for dense elimination; too many regions now drop the relative cut. Scores 0.28-0.32 under desktop load} \\
13 & 0.4321 & keep & \parbox[t]{9.2cm}{\raggedright Matrix parsing in C, one pass over the nonzeros; faces and blocks bit-identical to the numpy version} \\
29 & 0.538 & discard & \parbox[t]{9.2cm}{\raggedright Cycle V/W, K\_THRESH, ALPHA, REACTION, RHO, TAU, COARSEST, MIN\_REGION, ANISO: none beats base across the suite. K a.5 ties. Parametric ceiling; loop stopped} \\
30 & 0.565 & keep & \parbox[t]{9.2cm}{\raggedright Grouping probe: if the block-triangular solve contracts slower than 0.7 on the system aggregated to 8\^{}3 boxes, all coupled fields form one group. Reduction moved into flux\_mg.py; bit-identical to amfr.AMFR with the probe off} \\
31 & -- & discard & \parbox[t]{9.2cm}{\raggedright Gate region vectors by the non-Galerkin share of the coarse reaction (DEFL\_FRAC): separates per suite but a 3T group with share 0.03 still needs them (22 -\textgreater{} 39 its); per-region and probe criteria do not separate either} \\
32 & -- & discard & \parbox[t]{9.2cm}{\raggedright Finest-level region correction after the cycle only (one Q, one matvec)} \\
33 & 0.573 / 0.473 vs base 0.550 / 0.527 (paired) & discard & \parbox[t]{9.2cm}{\raggedright Balanced deflation through precomputed H = AW, G = A\^{}T W (no extra matvec), F = 1 sweep, KMIN\_RATIO 4. Patches in round3\_cost/. Setup, mostly build\_regions/assemble\_regions, is now the main cost} \\
\end{longtable}}
\FloatBarrier
\section{SSR（锁定版本 G，11 月 8 日）}
每行：SSR 迭代和时间；BoomerAMG \texttt{fast.25}；每个算例最优的 BoomerAMG 设置（含 AIR 和 systems 设置）。比值 = SSR / 对照，失败按 8 倍计。判据与 FluxMG 相同（双双精度原残差 $<10^{-8}$），3 次打乱顺序重复。来源 \texttt{gensolver/round3/final\_*\_G.log}。
{\scriptsize\setlength{\tabcolsep}{3pt}
\begin{longtable}{lrrrrrlrrr}
\caption{开发集（23）。对 fast.25 的 SCORE: 0.126  胜 11/23；对每例最优的 SCORE: -0.320  胜 8/23；SSR 失败: ['elast:33:0.3:1e4']}\\
\toprule
算例 & SSR & (s) & fast.25 & (s) & 比值 & 最优设置 & 迭代 & (s) & 比值 \\
\midrule
\endfirsthead
\toprule
算例 & SSR & (s) & fast.25 & (s) & 比值 & 最优设置 & 迭代 & (s) & 比值 \\
\midrule
\endhead
\bottomrule
\endlastfoot
\texttt{poisson3d:64} & 22 & 0.62 & 17 & 0.51 & 1.20 & \texttt{fast.25} & 17 & 0.51 & 1.20 \\
\texttt{aniso:64:100} & 44 & 0.53 & 16 & 0.59 & \textbf{0.91} & \texttt{fast.5} & 16 & 0.55 & \textbf{0.97} \\
\texttt{rot2d:512:1e-3:30} & 144 & 2.76 & 92 & 2.11 & 1.31 & \texttt{noagg.25} & 17 & 0.86 & 3.22 \\
\texttt{rot3d:48:1e-2:1e-4} & 62 & 0.69 & 48 & 0.54 & 1.28 & \texttt{fast.7} & 37 & 0.51 & 1.36 \\
\texttt{fe3d:150000:1:1} & 44 & 1.24 & 46 & 1.4 & \textbf{0.89} & \texttt{fast.9} & 39 & 1.19 & 1.04 \\
\texttt{fe3d:150000:1e6:2} & 49 & 1.33 & 46 & 1.39 & \textbf{0.95} & \texttt{fast.9} & 41 & 1.23 & 1.08 \\
\texttt{stretch:48:100} & 48 & 0.59 & 38 & 0.54 & 1.10 & \texttt{fast.5} & 31 & 0.43 & 1.37 \\
\texttt{lognorm3} & 25 & 0.66 & 19 & 0.56 & 1.18 & \texttt{fast.25} & 19 & 0.56 & 1.18 \\
\texttt{tubes} & 35 & 0.78 & 20 & 0.66 & 1.18 & \texttt{fast.25} & 20 & 0.66 & 1.18 \\
\texttt{perc3} & 27 & 0.65 & 19 & 0.68 & \textbf{0.95} & \texttt{fast.25} & 19 & 0.68 & \textbf{0.95} \\
\texttt{fourth:512} & 26 & 0.66 & 18 & 0.57 & 1.16 & \texttt{falgout.25} & 8 & 0.54 & 1.21 \\
\texttt{cd2d:512:1e-3:rec} & 18 & 0.6 & 34 & 0.68 & \textbf{0.88} & \texttt{noagg.25} & 20 & 0.62 & \textbf{0.97} \\
\texttt{cd2d:512:1e-5:rec} & 16 & 0.48 & 26 & 0.61 & \textbf{0.79} & \texttt{noagg.25} & 19 & 0.55 & \textbf{0.87} \\
\texttt{cd3d:64:1e-3} & 12 & 0.7 & 21 & 1.52 & \textbf{0.46} & \texttt{fast.9} & 39 & 0.9 & \textbf{0.78} \\
\texttt{elast:41:0.3:1} & 600 & 19.38 & 400 & 19.26 & 1.01 & \texttt{nodal.5} & 47 & 2.97 & 6.51 \\
\texttt{elast:41:0.45:1} & 1200 & 37.87 & 800 & 38.32 & \textbf{0.99} & \texttt{nodal.5} & 69 & 4.24 & 8.00 \\
\texttt{elast:33:0.3:1e4} & 1200 & 失败 & 1200 & 失败 & 8.00 & \texttt{nodal.5} & 1000 & 26.56 & 8.00 \\
\texttt{graph:ws:200000:1} & 21 & 0.57 & 24 & 1.43 & \textbf{0.40} & \texttt{fast.9} & 20 & 0.63 & \textbf{0.91} \\
\texttt{graph:ba:200000:1} & 14 & 0.53 & 10 & 3.74 & \textbf{0.14} & \texttt{fast.9} & 10 & 3.61 & \textbf{0.15} \\
\texttt{graph:rgg:200000:1} & 28 & 1.1 & 23 & 0.93 & 1.18 & \texttt{fast.25} & 23 & 0.93 & 1.18 \\
\texttt{g3t:64:2e-11:2e-4} & 25 & 2.54 & 30 & 2.27 & 1.12 & \texttt{fast.25} & 30 & 2.27 & 1.12 \\
\texttt{gen:64:3:2:1:0:2e-2:1} & 34 & 3.08 & 30 & 2.3 & 1.34 & \texttt{fast.25} & 30 & 2.3 & 1.34 \\
\texttt{mg:32:8:1e2:1e-2:1:7} & 22 & 0.47 & 109 & 2.36 & \textbf{0.20} & \texttt{fast.7} & 109 & 2.35 & \textbf{0.20} \\
\end{longtable}}
{\scriptsize\setlength{\tabcolsep}{3pt}
\begin{longtable}{lrrrrrlrrr}
\caption{留出集（14）。对 fast.25 的 SCORE: 0.192  胜 6/14；对每例最优的 SCORE: -0.447  胜 3/14；SSR 失败: []}\\
\toprule
算例 & SSR & (s) & fast.25 & (s) & 比值 & 最优设置 & 迭代 & (s) & 比值 \\
\midrule
\endfirsthead
\toprule
算例 & SSR & (s) & fast.25 & (s) & 比值 & 最优设置 & 迭代 & (s) & 比值 \\
\midrule
\endhead
\bottomrule
\endlastfoot
\texttt{poisson3d:100} & 25 & 2.63 & 19 & 2.16 & 1.22 & \texttt{fast.25} & 19 & 2.16 & 1.22 \\
\texttt{rot2d:700:1e-4:60} & 215 & 7.66 & 135 & 6.23 & 1.23 & \texttt{noagg.25} & 20 & 1.43 & 5.38 \\
\texttt{rot3d:60:1e-3:1e-3} & 83 & 1.83 & 63 & 1.52 & 1.21 & \texttt{fast.7} & 45 & 1.3 & 1.41 \\
\texttt{fe3d:300000:1e4:5} & 78 & 3.97 & 71 & 4.31 & \textbf{0.92} & \texttt{fast.9} & 55 & 3.46 & 1.15 \\
\texttt{stretch:60:1000} & 53 & 1.33 & 41 & 1.23 & 1.08 & \texttt{fast.7} & 37 & 1.03 & 1.30 \\
\texttt{cd2d:700:1e-4:rec} & 16 & 1.11 & 28 & 1.28 & \textbf{0.86} & \texttt{noagg.25} & 19 & 1.07 & 1.03 \\
\texttt{cd2d:512:1e-3:uni} & 14 & 0.53 & 26 & 0.53 & 1.01 & \texttt{noagg.25} & 13 & 0.4 & 1.32 \\
\texttt{cd3d:80:1e-4} & 14 & 1.4 & 21 & 3.16 & \textbf{0.44} & \texttt{fast.9} & 41 & 1.83 & \textbf{0.77} \\
\texttt{elast:51:0.4:1} & 1000 & 63.67 & 1000 & 97.45 & \textbf{0.65} & \texttt{nodal.5} & 55 & 6.94 & 8.00 \\
\texttt{graph:ws:300000:2} & 21 & 0.86 & 23 & 2.35 & \textbf{0.37} & \texttt{fast.9} & 24 & 1.14 & \textbf{0.76} \\
\texttt{graph:rgg:300000:2} & 31 & 1.75 & 25 & 1.52 & 1.15 & \texttt{fast.25} & 25 & 1.52 & 1.15 \\
\texttt{gen:64:10:2:1:0:2e-2:5} & 46 & 3.96 & 30 & 2.41 & 1.64 & \texttt{fast.25} & 30 & 2.41 & 1.64 \\
\texttt{mg:32:4:1e2:1:1:8} & 14 & 0.22 & 26 & 0.48 & \textbf{0.45} & \texttt{fast.25} & 26 & 0.48 & \textbf{0.45} \\
\texttt{insul} & 26 & 0.65 & 19 & 0.6 & 1.09 & \texttt{fast.7} & 21 & 0.63 & 1.03 \\
\end{longtable}}
{\scriptsize\setlength{\tabcolsep}{3pt}
\begin{longtable}{lrrrrrlrrr}
\caption{BoomerAMG 弱项集 WEAK（11）。对 fast.25 的 SCORE: 1.919  胜 10/11；对每例最优的 SCORE: 1.420  胜 9/11；SSR 失败: []}\\
\toprule
算例 & SSR & (s) & fast.25 & (s) & 比值 & 最优设置 & 迭代 & (s) & 比值 \\
\midrule
\endfirsthead
\toprule
算例 & SSR & (s) & fast.25 & (s) & 比值 & 最优设置 & 迭代 & (s) & 比值 \\
\midrule
\endhead
\bottomrule
\endlastfoot
\texttt{cdmix:512:1e-5:1} & 32 & 0.75 & 65 & 1.34 & \textbf{0.56} & \texttt{noagg.25} & 20 & 0.6 & 1.25 \\
\texttt{cdmix:512:1e-3:2} & 33 & 0.85 & 38 & 0.77 & 1.11 & \texttt{noagg.25} & 20 & 0.63 & 1.35 \\
\texttt{cd3m:64:1e-4:4} & 18 & 0.76 & 31 & 1.44 & \textbf{0.53} & \texttt{air.2} & 31 & 1.1 & \textbf{0.69} \\
\texttt{cd3d:64:1e-5} & 14 & 0.62 & 21 & 1.54 & \textbf{0.40} & \texttt{air.1} & 29 & 0.72 & \textbf{0.87} \\
\texttt{markov:ba:200000:0.99:1} & 18 & 0.87 & 69 & 1.93 & \textbf{0.45} & \texttt{fast.5} & 60 & 1.49 & \textbf{0.58} \\
\texttt{markov:ws:200000:0.999:1} & 25 & 1.04 & 271 & 12.59 & \textbf{0.08} & \texttt{fast.9} & 272 & 5.67 & \textbf{0.18} \\
\texttt{dlap:ba:200000:1} & 28 & 1.12 & 127 & 7.01 & \textbf{0.16} & \texttt{air.1} & 197 & 3.77 & \textbf{0.30} \\
\texttt{dlap:ws:200000:1} & 28 & 1.14 & 400 & 7.4 & \textbf{0.15} & \texttt{fast.7} & 400 & 7.29 & \textbf{0.16} \\
\texttt{mg:32:16:1e2:1e-2:1:3} & 23 & 0.82 & 127 & 5.28 & \textbf{0.15} & \texttt{sys.n} & 127 & 3.18 & \textbf{0.26} \\
\texttt{mg:32:8:1e3:1e-3:0:4} & 11 & 0.2 & 1200 & 21.86 & \textbf{0.01} & \texttt{sys.u} & 1200 & 14.42 & \textbf{0.01} \\
\texttt{graph:ba:300000:3} & 15 & 0.84 & 10 & 6.96 & \textbf{0.12} & \texttt{fast.9} & 10 & 6.75 & \textbf{0.12} \\
\end{longtable}}
{\scriptsize\setlength{\tabcolsep}{3pt}
\begin{longtable}{lrrrrrlrrr}
\caption{弱项留出 WEAK\_HOLD（6，曾用于找 bug）。对 fast.25 的 SCORE: 1.606  胜 6/6；对每例最优的 SCORE: 1.347  胜 5/6；SSR 失败: []}\\
\toprule
算例 & SSR & (s) & fast.25 & (s) & 比值 & 最优设置 & 迭代 & (s) & 比值 \\
\midrule
\endfirsthead
\toprule
算例 & SSR & (s) & fast.25 & (s) & 比值 & 最优设置 & 迭代 & (s) & 比值 \\
\midrule
\endhead
\bottomrule
\endlastfoot
\texttt{cdmix:700:1e-4:7} & 38 & 1.75 & 43 & 1.76 & \textbf{0.99} & \texttt{noagg.25} & 20 & 1.16 & 1.51 \\
\texttt{cd3m:72:1e-5:9} & 18 & 1.06 & 34 & 2.14 & \textbf{0.49} & \texttt{air.2} & 34 & 1.65 & \textbf{0.64} \\
\texttt{markov:ba:300000:0.995:5} & 23 & 1.48 & 112 & 4.64 & \textbf{0.32} & \texttt{air.1} & 134 & 3.45 & \textbf{0.43} \\
\texttt{dlap:ws:300000:6} & 113 & 5.56 & 398 & 11.4 & \textbf{0.49} & \texttt{fast.9} & 382 & 10.44 & \textbf{0.53} \\
\texttt{mg:32:12:3e2:3e-3:1:11} & 24 & 0.57 & 400 & 12.0 & \textbf{0.05} & \texttt{sys.u} & 315 & 7.06 & \textbf{0.08} \\
\texttt{graph:ba:250000:8} & 15 & 0.67 & 10 & 5.13 & \textbf{0.13} & \texttt{fast.7} & 10 & 5.03 & \textbf{0.13} \\
\end{longtable}}
{\scriptsize\setlength{\tabcolsep}{3pt}
\begin{longtable}{lrrrrrlrrr}
\caption{弱项干净留出 WEAK\_HOLD2（9）。对 fast.25 的 SCORE: 2.133  胜 9/9；对每例最优的 SCORE: 1.764  胜 8/9；SSR 失败: []}\\
\toprule
算例 & SSR & (s) & fast.25 & (s) & 比值 & 最优设置 & 迭代 & (s) & 比值 \\
\midrule
\endfirsthead
\toprule
算例 & SSR & (s) & fast.25 & (s) & 比值 & 最优设置 & 迭代 & (s) & 比值 \\
\midrule
\endhead
\bottomrule
\endlastfoot
\texttt{cdmix:600:1e-5:13} & 39 & 1.22 & 58 & 1.7 & \textbf{0.72} & \texttt{noagg.25} & 21 & 0.88 & 1.38 \\
\texttt{cd3m:60:1e-3:21} & 17 & 0.67 & 26 & 1.1 & \textbf{0.61} & \texttt{air.2} & 30 & 0.88 & \textbf{0.77} \\
\texttt{markov:ws:250000:0.9999:17} & 31 & 1.58 & 600 & 36.88 & \textbf{0.04} & \texttt{fast.9} & 600 & 15.87 & \textbf{0.10} \\
\texttt{markov:rgg:200000:0.99:3} & 17 & 1.25 & 66 & 1.9 & \textbf{0.66} & \texttt{fast.9} & 72 & 1.86 & \textbf{0.67} \\
\texttt{dlap:ba:300000:19} & 27 & 1.64 & 149 & 13.8 & \textbf{0.12} & \texttt{fast.9} & 149 & 5.63 & \textbf{0.29} \\
\texttt{dlap:rgg:200000:23} & 28 & 1.69 & 400 & 8.97 & \textbf{0.19} & \texttt{fast.9} & 400 & 8.74 & \textbf{0.19} \\
\texttt{mg:32:24:1e2:1e-2:1:29} & 23 & 1.17 & 152 & 9.55 & \textbf{0.12} & \texttt{sys.n} & 152 & 5.52 & \textbf{0.21} \\
\texttt{mg:24:8:3e3:3e-3:0:31} & 5 & 0.06 & 600 & 4.41 & \textbf{0.01} & \texttt{sys.n} & 1000 & 2.95 & \textbf{0.02} \\
\texttt{graph:ba:400000:37} & 14 & 1.19 & 10 & 10.62 & \textbf{0.11} & \texttt{fast.9} & 10 & 10.15 & \textbf{0.12} \\
\end{longtable}}


\end{document}
